% =====================================================================
% Supplementary Information for
% "Indirect reciprocity with inertial reputation updating"
% Scope: private assessment; discrete inertial (direct-reset) record.
% Compile with: pdflatex Supplementary_Information.tex (twice).
% Notation identical to Methods_body.tex (Supplementary Note 1).
% =====================================================================
\documentclass[fleqn,10pt]{article}
\usepackage[utf8]{inputenc}
\usepackage[T1]{fontenc}
\usepackage[english]{babel}
\usepackage[letterpaper,left=2.2cm,right=2.2cm,top=2.3cm,bottom=2.3cm]{geometry}
\usepackage{times,mathptmx}
\usepackage{amsmath,amssymb,amsthm,bm,booktabs,longtable,array,microtype}
\usepackage{titlesec}
\usepackage[colorlinks=true,allcolors=blue]{hyperref}
\usepackage{caption}
\captionsetup{font=small,labelfont=bf,labelsep=period,justification=raggedright,singlelinecheck=false}
\titleformat{\section}{\large\sffamily\bfseries}{Supplementary Note \thesection.}{0.5em}{}
\titleformat{\subsection}{\sffamily\bfseries}{\thesection.\arabic{subsection}}{0.5em}{}
\setlength{\parindent}{0pt}
\setlength{\parskip}{0.45em}
\setlength{\emergencystretch}{2em}
\allowdisplaybreaks
\renewcommand{\theequation}{S\arabic{equation}}
\renewcommand{\thetable}{S\arabic{table}}
\renewcommand{\tablename}{Supplementary Table}
\theoremstyle{plain}
\newtheorem{proposition}{Proposition}
\newtheorem{lemma}[proposition]{Lemma}
\newtheorem{corollary}[proposition]{Corollary}
\theoremstyle{remark}
\newtheorem{remark}[proposition]{Remark}
\renewcommand{\theproposition}{S\arabic{proposition}}
% ---- main-text figure numbers (Manuscript_formal_20261028.tex) ----
\newcommand{\FigModel}{1}
\newcommand{\FigRecognition}{2}
\newcommand{\FigStrategy}{3}
\newcommand{\FigDepth}{4}
\newcommand{\FigOpportunism}{5}
\newcommand{\EDRecognition}{1}
\newcommand{\EDResident}{2}
\newcommand{\EDPlanes}{3}
\newcommand{\EDEdge}{4}
\newcommand{\EDRatio}{5}
\newcommand{\EDErrors}{6}
\newcommand{\EDDepth}{7}
\newcommand{\EDReadout}{8}
\newcommand{\EDSecurity}{9}
% ---- short-hands (pure abbreviations; no new symbols) ----
\newcommand{\Rm}{\mathcal R}
\newcommand{\G}{\mathrm G}
\newcommand{\B}{\mathrm B}
\newcommand{\W}{\mathrm W}
\renewcommand{\S}{\mathrm S}
\newcommand{\C}{\mathrm C}
\newcommand{\D}{\mathrm D}
\newcommand{\ALLC}{\mathrm{ALLC}}
\newcommand{\ALLD}{\mathrm{ALLD}}
\newcommand{\DISC}{\mathrm{DISC}}
\newcommand{\Pp}[2]{P^+_{\mathrm{#1}\mid\mathrm{#2}}}
\newcommand{\rR}{r^*_{\mathrm R}}
\newcommand{\uR}{u^*_{\mathrm R}}

\begin{document}
\raggedbottom
\begin{center}
{\Large\sffamily\bfseries Supplementary Information}\\[0.6em]
{\large\sffamily Indirect reciprocity with inertial reputation updating}\\[0.4em]
{\small Wenqiang Zhu \textit{et al.}}
\end{center}

\vspace{0.5em}
\noindent This Supplementary Information gives the full derivations behind the
Methods. It treats private assessment and the discrete inertial
(direct-reset) reputation record used throughout the paper. Supplementary
Notes~\ref{sec:notation}--\ref{sec:strategy} derive the model and its
reputation dynamics, Notes~\ref{sec:invasion}--\ref{sec:filters} contain the
invasion analyses, and Notes~\ref{sec:simulations}--\ref{sec:limits} describe
simulations, numerical methods, the settings of every figure and the scope of
the results. Figure numbers refer to the main text and Extended Data.

{\small\tableofcontents}
\newpage

% =====================================================================
\section{Notation and parameters}
\label{sec:notation}

Supplementary Table~\ref{tab:symbols} lists every symbol used in the Methods
and in this Supplementary Information. Each symbol has a single meaning.
Two conventions should be noted.
(i) The letters $\G$ and $\B$ name both the two values of a binary image
(good, bad) and the two established latent states. Because an established state
is always read as the image of the same name
($\sigma(\G)=1$, $\sigma(\B)=0$), this use is unambiguous: latent states appear
only as arguments of $\sigma(\cdot)$, $\pi_\cdot$ and the transition rules,
whereas images appear only as conditions in $P^+_{A\mid Y}$ and in the
pair probabilities $h_{\G\G},h_{\G\B},h_{\B\B}$.
(ii) Subscripts $\C$ and $\D$ denote intended actions. An individual that always
intends $\C$ (ALLC) or $\D$ (ALLD) is assessed positively with probability
$\Phi_\C$ or $\Phi_\D$; strategy-level quantities carry full strategy names
($r_{\ALLC}$, $u_{\DISC}$, \dots). The subscript $\mathrm R$ denotes the
resident of a homogeneous DISC population and $\mathrm X$ a rare mutant.

\begin{longtable}{@{}p{0.27\linewidth}p{0.69\linewidth}@{}}
\caption{\textbf{Symbols.}}\label{tab:symbols}\\
\toprule
Symbol & Meaning \\
\midrule
\endfirsthead
\toprule
Symbol & Meaning (continued)\\
\midrule
\endhead
\bottomrule
\endfoot
\multicolumn{2}{@{}l}{\textit{Game, strategies and norms}}\\
$b,\ c$ & Benefit to the recipient and cost to the donor of realized cooperation; $b>c>0$.\\
$I,J\in\{\ALLC,\ALLD,\DISC\}$ & Strategies; $f_I$ frequency, $\boldsymbol f$ composition, $\sum_If_I=1$.\\
$A\in\{\C,\D\}$ & Intended action (cooperate, defect).\\
$Y\in\{\G,\B\}$ & Binary image of a recipient held by an observer (good, bad).\\
$\nu_{AY}$ & Probability that the norm prescribes a positive assessment of realized action $A$ towards a recipient with image $Y$: $\nu_{\C\G}=1$, $\nu_{\D\G}=0$, $\nu_{\C\B}=p$, $\nu_{\D\B}=q$.\\
$p,\ q$ & Norm parameters; SC $(1,0)$, SH $(0,0)$, SJ $(0,1)$, SS $(1,1)$.\\
$e_1,\ e_2$ & Implementation error (intended C becomes D) and assessment error (recorded assessment reversed).\\
$\epsilon,\ \kappa$ & $\epsilon=(1-e_1)(1-e_2)+e_1e_2$; $\kappa=(1-e_1)(1-2e_2)=\epsilon-e_2$.\\
$P^+_{A\mid Y}$ & Probability of a positive recorded assessment given intended action $A$ and the observer's image $Y$ of the recipient (errors included).\\
$\Phi_\C(r),\ \Phi_\D(r)$ & Probability that intended C or D is assessed positively when the observer regards the recipient as good with probability $r$.\\
$\Delta\Phi(r)$ & Assessment contrast $\Phi_\C(r)-\Phi_\D(r)=\kappa[r+(p-q)(1-r)]$.\\[0.3em]
\multicolumn{2}{@{}l}{\textit{Inertial record}}\\
$L_{ij}\in\mathcal L$ & Latent record of observer $i$ about target $j$; $\mathcal L=\{\G,\W_1,\dots,\W_{k_1},\B,\S_1,\dots,\S_{k_2}\}$.\\
$\W_m,\ \S_n$ & Warning states ($m=1,\dots,k_1$) and recovery states ($n=1,\dots,k_2$).\\
$k_1,\ k_2,\ K$ & Tolerance depth, strictness depth, total depth $K=k_1+k_2$.\\
$\alpha,\ \beta$ & Tolerance strength and strictness strength (readout parameters), $0\le\alpha,\beta\le1$.\\
$\sigma(\ell)$, $\boldsymbol\sigma$ & Probability that state $\ell$ is read as good; column vector of these probabilities.\\
$u,\ v$ & Probability that a recorded assessment of a target is positive; $v=1-u$.\\
$T^{+},T^{-},T(u)$ & Transition matrices after a positive or negative assessment; $T(u)=uT^{+}+vT^{-}$.\\
$\pi_\ell(u),\ Z(u)$ & Stationary distribution of a record with i.i.d.\ input $u$; normalization $Z=v^{k_1}+u^{k_2}-v^{k_1}u^{k_2}$.\\
$\Rm(u)$ & Reputation map $\Rm(u;\alpha,\beta,k_1,k_2)=\sum_\ell\pi_\ell(u)\sigma(\ell)$; $\Rm'=\mathrm d\Rm/\mathrm du$.\\
$\Sigma_\W,\ \Sigma_\S,\ \Omega$ & Warning-path weight, recovery-path weight and numerator of $\Rm$ (Supplementary Note~\ref{sec:map}).\\
$\Theta_{\G\to\B},\ \Theta_{\B\to\G}$ & Number of assessments needed to leave good (bad) standing, starting from $\G$ ($\B$).\\[0.3em]
\multicolumn{2}{@{}l}{\textit{Population reputations}}\\
$r_I,\ \bar r$ & Probability that an observer regards a type-$I$ individual as good; $\bar r=\sum_If_Ir_I$.\\
$h_{\G\G},h_{\G\B},h_{\B\B}$ & Probabilities that donor and observer both regard a random recipient as good, hold a given ordered disagreement, or both regard it as bad.\\
$u_I$ & Probability that a type-$I$ donor is assessed positively.\\
$\boldsymbol\pi_I$ & Row vector of latent-state probabilities of records about type-$I$ targets.\\
$\tau,\ t',\ t$ & Auxiliary relaxation time; time in assessments per record; evolutionary time.\\
$u_{\mathrm R}(r),\ F(r)$ & Resident assessment probability $r\Phi_\C(r)+(1-r)\Phi_\D(r)$ in a homogeneous DISC population; $F(r)=\Rm[u_{\mathrm R}(r)]-r$.\\
$\rR,\ \uR,\ \Phi_A^*,\ \Delta\Phi^*$ & Selected resident reputation (highest strictly stable root); $u_{\mathrm R}(\rR)$; $\Phi_A(\rR)$; $\Delta\Phi(\rR)$.\\[0.3em]
\multicolumn{2}{@{}l}{\textit{Evolution and cooperation}}\\
$\Pi_I,\ \bar\Pi$ & Expected payoff of strategy $I$; population mean payoff.\\
$\rho_{\mathrm{int}},\ \rho$ & Intended and realized cooperation rate, $\rho=(1-e_1)\rho_{\mathrm{int}}$.\\
$\bar\rho,\ \omega_\C,\ \bar\rho_\C$ & Mean long-run cooperation over initial compositions; cooperative-basin size; cooperation within the basin.\\
$n_{\mathrm{init}}$ & Number of initial compositions (512 per sample).\\[0.3em]
\multicolumn{2}{@{}l}{\textit{Invasion}}\\
$\chi_{\mathrm X}$ & Probability that a rare mutant X intends to help a resident recipient.\\
$u_{\mathrm X},\ r^*_{\mathrm X}$ & Positive-assessment probability and stationary reputation of the mutant.\\
$\Lambda_{\mathrm X}$ & Normalized payoff advantage $(\Pi_{\mathrm X}-\Pi_{\mathrm R})/[c(1-e_1)]$.\\
$\lambda_{\ALLD},\ \lambda_{\ALLC}$ & Unnormalized payoff advantages of rare ALLD and ALLC over residents.\\
$r^*_{\ALLD},\ r^*_{\ALLC}$ & Reputations of rare ALLD and ALLC in the resident population, $\Rm(\Phi^*_\D)$ and $\Rm(\Phi^*_\C)$.\\
$R_{\mathrm L},\ R_{\mathrm U}$ & Lower and upper benefit-to-cost thresholds for resistance to ALLD and ALLC.\\
$\delta$ & Probability that O-DISC withholds an intended cooperation.\\
$\Lambda(\delta),\ H(\delta)$ & Normalized O-DISC advantage; reputation lost per unit of withheld cooperation.\\
$H_{\min},\ \delta_\mathrm c$ & $\min_{0\le\delta\le1}H(\delta)$; upper end of the interval with $\Lambda>0$.\\
$\alpha_{\mathrm L},\ \alpha_{\mathrm U}$ & Lower (fold) and upper (invasion) boundaries of a secure readout interval.\\
$\Rm_{\mathrm{agg}},\ M,\ Q$ & Aggregation map; number of sampled assessments; required number of positive ones.\\
$\Rm_{\mathrm{score}},\ z,\ z_{\max},\ z_0$ & Quantitative-score map; score; score bound; good-image threshold.\\[0.3em]
\multicolumn{2}{@{}l}{\textit{Finite-population simulations}}\\
$N,\ n_I$ & Population size and number of type-$I$ individuals.\\
$s_{ij}$ & $\sigma(L_{ij})$, the probability that DISC observer $i$ currently reads target $j$ as good.\\
$g_j,\ g_I$ & Mean goodness of target $j$; average over type-$I$ targets and sampling times.\\
$d_j,\ d_I$ & Probability that two distinct DISC observers read target $j$ differently; average over type-$I$ targets and times.\\
$\Delta g,\ \bar d$ & Reputational separation and mean disagreement.\\
\end{longtable}

\begin{table}[htbp]
\centering\small
\caption{\textbf{Reference parameters and inertia settings.} Values apply unless a figure axis varies them.}
\label{tab:parameters}
\begin{tabular}{@{}lll@{}}
\toprule
Quantity & Value & Varied in\\
\midrule
Benefit, cost & $b=5$, $c=1$ & $b/c$: Extended Data Figs.~\EDRatio, \EDSecurity\\
Errors & $e_1=e_2=0.02$ & $e_1,e_2\in[0,0.45]$: Extended Data Fig.~\EDErrors\\
No inertia & $(k_1,k_2)=(0,0)$; equivalently $\alpha=\beta=0$ & --\\
Full one-step inertia & $(k_1,k_2)=(1,1)$, $\alpha=\beta=1$ & $(\alpha,\beta)$: Extended Data Figs.~\EDPlanes, \EDReadout\\
Selected multi-step inertia & SC $(2,0)$, SH $(3,0)$, SJ $(4,3)$, SS $(1,2)$; $\alpha=\beta=1$ & $\alpha$: Fig.~\FigOpportunism, Extended Data Fig.~\EDSecurity\\
Depth scans & $\alpha=\beta=1$ & $k_1,k_2$: Figs.~\FigRecognition, \FigDepth; Extended Data Figs.~\EDRecognition, \EDPlanes, \EDDepth\\
Initial records & all latent records in $\G$ ($r_I=1$) & all-bad: Extended Data Fig.~\EDResident e\\
\bottomrule
\end{tabular}
\end{table}

% =====================================================================
\section{Interactions, social norms and errors}
\label{sec:assessment}

\subsection{Sequence of an interaction}
A donor, a recipient and an observer are drawn from the population, the
observer independently of the two others. Each interaction proceeds as follows.
\begin{enumerate}
\item The donor forms an intended action. ALLC intends $\C$, ALLD intends $\D$,
and DISC consults its own record of the recipient, samples a binary image and
intends $\C$ if and only if the image is $\G$.
\item Intended cooperation is implemented as defection with probability $e_1$;
intended defection is always implemented. The recipient receives $b$ and the
donor pays $c$ only if cooperation is realized.
\item The observer consults its own record of the recipient, samples a binary
image $Y$ and applies the norm to the realized action: a positive assessment is
prescribed with probability $\nu_{AY}$.
\item The prescribed assessment is reversed with probability $e_2$.
\item The recorded assessment updates the observer's record of the donor. No
other record changes.
\end{enumerate}
The donor's and the observer's images of the recipient come from different
records, so the donor may act on a different view of the recipient from the
one used to judge it. This is the defining feature of private assessment.

\subsection{Error-integrated assessment probabilities}
Conditional on a realized action $A$ and an image $Y$, the recorded assessment
is positive with probability $\nu_{AY}(1-e_2)+(1-\nu_{AY})e_2=e_2+(1-2e_2)\nu_{AY}$.
Intended cooperation is realized as $\C$ with probability $1-e_1$ and as $\D$
with probability $e_1$, whereas intended defection is always realized as $\D$.
Averaging over the realized action,
\begin{equation}
P^+_{\C\mid Y}=e_2+(1-2e_2)\bigl[(1-e_1)\nu_{\C Y}+e_1\nu_{\D Y}\bigr],\qquad
P^+_{\D\mid Y}=e_2+(1-2e_2)\nu_{\D Y}.
\label{eq:S-assessment-general}
\end{equation}
With $\nu_{\C\G}=1$, $\nu_{\D\G}=0$, $\nu_{\C\B}=p$ and $\nu_{\D\B}=q$,
\begin{equation}
\Pp CG=\epsilon,\qquad \Pp DG=e_2,\qquad
\Pp CB=e_2+(1-2e_2)\bigl[(1-e_1)p+e_1q\bigr],\qquad
\Pp DB=e_2+(1-2e_2)q,
\label{eq:S-assessment}
\end{equation}
where $\epsilon=e_2+\kappa=(1-e_1)(1-e_2)+e_1e_2$ and
$\kappa=(1-e_1)(1-2e_2)$. Equivalently,
$\Pp CB=p(\epsilon-e_2)+q(1-\epsilon-e_2)+e_2$, because
$\epsilon-e_2=(1-e_1)(1-2e_2)$ and $1-\epsilon-e_2=e_1(1-2e_2)$. Supplementary
Table~\ref{tab:norms} gives the four norms. These probabilities already include
implementation error; the separate factor $1-e_1$ in the payoffs accounts for
realized donations, not for assessments.

\begin{table}[htbp]
\centering\small
\caption{\textbf{Norms and error-integrated assessment probabilities.}}
\label{tab:norms}
\begin{tabular}{@{}lcccccc@{}}
\toprule
Norm & $p=\nu_{\C\B}$ & $q=\nu_{\D\B}$ & $\Pp CG$ & $\Pp DG$ & $\Pp CB$ & $\Pp DB$\\
\midrule
Scoring (SC) & 1 & 0 & $\epsilon$ & $e_2$ & $\epsilon$ & $e_2$\\
Shunning (SH) & 0 & 0 & $\epsilon$ & $e_2$ & $e_2$ & $e_2$\\
Stern Judging (SJ) & 0 & 1 & $\epsilon$ & $e_2$ & $1-\epsilon$ & $1-e_2$\\
Simple Standing (SS) & 1 & 1 & $\epsilon$ & $e_2$ & $1-e_2$ & $1-e_2$\\
\bottomrule
\end{tabular}
\end{table}

\subsection{Action-conditional assessment probabilities}
Suppose that the observer regards the recipient as good with probability $r$,
independently of the donor's image. Then intended cooperation and intended
defection are assessed positively with probabilities
\begin{equation}
\Phi_\C(r)=r\Pp CG+(1-r)\Pp CB,\qquad
\Phi_\D(r)=r\Pp DG+(1-r)\Pp DB .
\label{eq:S-Phi}
\end{equation}
Both are linear in $r$, with slopes $\Phi_\C'=\Pp CG-\Pp CB$ and
$\Phi_\D'=\Pp DG-\Pp DB$. Their difference is
\begin{equation}
\Delta\Phi(r)=r(\Pp CG-\Pp DG)+(1-r)(\Pp CB-\Pp DB)=\kappa\,[r+(p-q)(1-r)],
\label{eq:S-contrast}
\end{equation}
because $\Pp CG-\Pp DG=\kappa$ and $\Pp CB-\Pp DB=(p-q)\kappa$. Thus
$\Delta\Phi=\kappa$ (SC), $\kappa r$ (SH, SS) and $\kappa(2r-1)$ (SJ): under SJ the
contrast changes sign at $r=1/2$, so on a low-reputation branch defection is
assessed better than cooperation.

A DISC donor in a homogeneous DISC population with reputation $r$ regards the
recipient as good with probability $r$, independently of the observer
(Supplementary Note~\ref{sec:population}). It therefore intends $\C$ with
probability $r$ and is assessed positively with probability
\begin{equation}
u_{\mathrm R}(r)=r\,\Phi_\C(r)+(1-r)\,\Phi_\D(r)
=r^2\Pp CG+r(1-r)\bigl(\Pp CB+\Pp DG\bigr)+(1-r)^2\Pp DB,
\label{eq:S-uR}
\end{equation}
with derivative $u_{\mathrm R}'(r)=\Delta\Phi(r)+r\Phi_\C'+(1-r)\Phi_\D'$.
Supplementary Table~\ref{tab:resident-functions} lists these functions.

\begin{table}[htbp]
\centering\small
\caption{\textbf{Resident assessment functions.} $\kappa=(1-e_1)(1-2e_2)$. For $e_1=e_2=0.02$: $\epsilon=0.9608$, $\kappa=0.9408$.}
\label{tab:resident-functions}
\begin{tabular}{@{}lllll@{}}
\toprule
Norm & $\Phi_\C(r)$ & $\Phi_\D(r)$ & $\Delta\Phi(r)$ & $u_{\mathrm R}(r)$\\
\midrule
SC & $\epsilon$ & $e_2$ & $\kappa$ & $e_2+\kappa r$\\
SH & $e_2+\kappa r$ & $e_2$ & $\kappa r$ & $e_2+\kappa r^2$\\
SJ & $1-\epsilon-(1-2\epsilon)r$ & $1-e_2-(1-2e_2)r$ & $\kappa(2r-1)$ & $1-e_2-(1-2e_2+\kappa)r+2\kappa r^2$\\
SS & $1-e_2-(1-e_2-\epsilon)r$ & $1-e_2-(1-2e_2)r$ & $\kappa r$ & $1-e_2-(1-2e_2)r+\kappa r^2$\\
\bottomrule
\end{tabular}
\end{table}

% =====================================================================
\section{The inertial reputation record}
\label{sec:record}

\subsection{States and transitions}
Each observer holds, for each target, a latent record in
$\mathcal L=\{\G,\W_1,\dots,\W_{k_1},\B,\S_1,\dots,\S_{k_2}\}$ with
$k_1,k_2\in\{0,1,2,\dots\}$. The established states $\G$ and $\B$ represent
good and bad standing. Warning states $\W_m$ count consecutive negative
assessments that challenge good standing, and recovery states $\S_n$ count
consecutive positive assessments that challenge bad standing. Supplementary
Table~\ref{tab:transitions} lists the transitions. Two rules generate it:
evidence against the established judgement advances along the corresponding
path, and evidence supporting it resets the whole path (direct reset). The
record has $K+2$ states. Consecutiveness refers to the assessments made by one
observer about one target; these need not be consecutive interactions of the
target.

\begin{table}[htbp]
\centering\small
\caption{\textbf{Transitions of a latent record.} Rows for $\W_m$ and $\S_n$ are absent when $k_1=0$ or $k_2=0$, respectively.}
\label{tab:transitions}
\begin{tabular}{@{}lll@{}}
\toprule
Current state & After a positive assessment & After a negative assessment\\
\midrule
$\G$ & $\G$ & $\W_1$ if $k_1\ge1$; $\B$ if $k_1=0$\\
$\W_m$, $1\le m<k_1$ & $\G$ & $\W_{m+1}$\\
$\W_{k_1}$ & $\G$ & $\B$\\
$\B$ & $\S_1$ if $k_2\ge1$; $\G$ if $k_2=0$ & $\B$\\
$\S_n$, $1\le n<k_2$ & $\S_{n+1}$ & $\B$\\
$\S_{k_2}$ & $\G$ & $\B$\\
\bottomrule
\end{tabular}
\end{table}

We order the states as $(\G,\W_1,\dots,\W_{k_1},\B,\S_1,\dots,\S_{k_2})$ and
write $T^{+}$ and $T^{-}$ for the $(K+2)\times(K+2)$ matrices with entry
$[T^\pm]_{\ell\ell'}=1$ if a positive (negative) assessment moves $\ell$ to
$\ell'$ and $0$ otherwise. Each row of $T^\pm$ contains exactly one 1. If a record
receives a positive assessment with probability $u$ and a negative one with
probability $v=1-u$, its state follows the Markov chain with transition matrix
\begin{equation}
T(u)=u\,T^{+}+v\,T^{-}.
\label{eq:S-T}
\end{equation}
For one-step inertia, in the order $(\G,\W_1,\B,\S_1)$,
\[
T(u)=\begin{pmatrix}u&v&0&0\\u&0&v&0\\0&0&v&u\\u&0&v&0\end{pmatrix}.
\]

\subsection{Binary readout}
When a record in state $\ell$ is consulted, it yields the image $\G$ with
probability $\sigma(\ell)$ and $\B$ otherwise, where
\begin{equation}
\sigma(\G)=1,\qquad\sigma(\W_m)=\alpha^m,\qquad\sigma(\B)=0,\qquad
\sigma(\S_n)=1-\beta^n,\qquad 0\le\alpha,\beta\le1 .
\label{eq:S-readout}
\end{equation}
Each consultation draws a fresh image, independently of all other
consultations given the current state, and does not change the state. In
particular, when a DISC donor is also the observer of its own action, the image
used to act and the image used to judge are independent draws from the same
record. The geometric form makes successive warning states progressively less
likely to be treated as good, and successive recovery states progressively more
likely to be treated as good. The depths control how much consistent evidence
completes a reversal of the latent state; the readout parameters control how
strongly incomplete evidence already affects behaviour and assessment.

\subsection{Settings that reduce to instantaneous updating}
\begin{lemma}\label{lem:instantaneous}
If $(k_1,k_2)=(0,0)$, or if $\alpha=\beta=0$ at any depths, the image of a record
that has received at least one assessment equals its most recent assessment,
and a record that has received none yields $\G$.
\end{lemma}
\begin{proof}
For $(k_1,k_2)=(0,0)$ the record has the states $\G$ and $\B$ only, and every
assessment sets it to $\G$ or $\B$ according to its sign. For $\alpha=\beta=0$
and any depths, the states that can follow a positive assessment are $\G$ and
$\S_1,\dots,\S_{k_2}$, all with $\sigma=1$ because $1-0^n=1$ for $n\ge1$; the
states that can follow a negative assessment are $\W_1,\dots,\W_{k_1}$ and $\B$,
all with $\sigma=0$ because $0^m=0$ for $m\ge1$. Initially every record is in
$\G$.
\end{proof}
The two implementations therefore generate identical image processes, not
only identical stationary statistics. We use $(0,0)$ in the depth scans and in
the stationary analyses, and the one-step state space with $\alpha=\beta=0$ in
the instantaneous panels of Figs.~\FigRecognition a--d and~\FigStrategy.

% =====================================================================
\section{The stationary reputation map}
\label{sec:map}

\subsection{Stationary distribution}
\begin{proposition}[Stationary distribution]\label{prop:stationary}
For $0<u<1$ the chain $T(u)$ is irreducible and aperiodic, and its unique
stationary distribution is
\begin{equation}
\pi_\G=\frac{u^{k_2+1}}{Z},\qquad
\pi_{\W_m}=v^m\pi_\G,\qquad
\pi_\B=\frac{v^{k_1+1}}{Z},\qquad
\pi_{\S_n}=u^n\pi_\B,\qquad
Z(u)=v^{k_1}+u^{k_2}-v^{k_1}u^{k_2},
\label{eq:S-stationary}
\end{equation}
for $1\le m\le k_1$ and $1\le n\le k_2$.
\end{proposition}
\begin{proof}
Irreducibility: from any state, a run of negative assessments reaches $\B$ and a
run of positive assessments then reaches $\G$; from $\G$ every warning state and
from $\B$ every recovery state is reachable. Aperiodicity: $\G$ has a self-loop
with probability $u>0$. Hence a unique stationary distribution exists.

\textit{Paths.} The state $\W_m$ ($m\ge1$) is entered only from $\W_{m-1}$, with
$\W_0\equiv\G$, after a negative assessment. Stationarity of $\W_m$ requires
$\pi_{\W_m}=v\,\pi_{\W_{m-1}}$, hence $\pi_{\W_m}=v^m\pi_\G$. Likewise, $\S_n$ is
entered only from $\S_{n-1}$, with $\S_0\equiv\B$, after a positive assessment,
so $\pi_{\S_n}=u^n\pi_\B$.

\textit{Flux balance.} Partition $\mathcal L$ into the good-standing block
$\{\G,\W_1,\dots,\W_{k_1}\}$ and the bad-standing block
$\{\B,\S_1,\dots,\S_{k_2}\}$. The only transition from the first block to the
second is the negative assessment from $\W_{k_1}$ (from $\G$ if $k_1=0$), with
stationary flux $v\,\pi_{\W_{k_1}}=v^{k_1+1}\pi_\G$. The only transition in the
opposite direction is the positive assessment from $\S_{k_2}$ (from $\B$ if
$k_2=0$), with flux $u^{k_2+1}\pi_\B$. In a stationary state the two fluxes are
equal:
\begin{equation}
v^{k_1+1}\pi_\G=u^{k_2+1}\pi_\B .
\label{eq:S-flux}
\end{equation}
\textit{Normalization.} Writing $\pi_\G=u^{k_2+1}/Z$ and $\pi_\B=v^{k_1+1}/Z$
satisfies Eq.~\eqref{eq:S-flux}. Summing all probabilities,
\[
Z=u^{k_2+1}\sum_{m=0}^{k_1}v^m+v^{k_1+1}\sum_{n=0}^{k_2}u^n
=u^{k_2}\bigl(1-v^{k_1+1}\bigr)+v^{k_1}\bigl(1-u^{k_2+1}\bigr)
=v^{k_1}+u^{k_2}-v^{k_1}u^{k_2}(u+v),
\]
which is the stated $Z$ because $u+v=1$.
\end{proof}

The exponents record the evidence requirements: completing a reversal of good
standing requires $k_1+1$ negative assessments in a row, which suppresses
$\pi_\B$ by $v^{k_1+1}$, and completing a recovery requires $k_2+1$ positive
assessments in a row, which suppresses $\pi_\G$ by $u^{k_2+1}$.

\subsection{The reputation map}
Averaging the readout probabilities over Eq.~\eqref{eq:S-stationary} gives the
probability that a target is regarded as good in the stationary state of a
record with i.i.d.\ input $u$:
\begin{equation}
\Rm(u)\equiv\Rm(u;\alpha,\beta,k_1,k_2)
=\pi_\G+\sum_{m=1}^{k_1}\alpha^m\pi_{\W_m}+\sum_{n=1}^{k_2}(1-\beta^n)\pi_{\S_n}
=\frac{u^{k_2+1}\,\Sigma_\W(v)+v^{k_1+1}\,\Sigma_\S(u)}{Z(u)},
\label{eq:S-map}
\end{equation}
\[
\Sigma_\W(v)=\sum_{m=0}^{k_1}(\alpha v)^m,\qquad
\Sigma_\S(u)=\sum_{n=1}^{k_2}u^n(1-\beta^n).
\]
Empty sums vanish. At the endpoints, the chain is absorbed in $\B$ ($u=0$) or
$\G$ ($u=1$), so $\Rm(0)=0$ and $\Rm(1)=1$; Eq.~\eqref{eq:S-map} attains these
values by continuity (for $k_1=0$ or $k_2=0$ it is already regular there).
$\Rm(u)$ is exact for a record with i.i.d.\ input; how $u$ is determined in a
population is the subject of Supplementary Note~\ref{sec:population}.

\textit{Special cases.}
\begin{enumerate}
\item No inertia. For $(k_1,k_2)=(0,0)$, $Z=1$ and $\Rm(u)=u$. By
Lemma~\ref{lem:instantaneous}, the same holds for $\alpha=\beta=0$ at any depths.
\item One-step inertia, $(k_1,k_2)=(1,1)$:
\begin{equation}
\Rm(u)=\frac{u^2\,[1+\alpha(1-u)]+(1-\beta)\,u(1-u)^2}{1-u+u^2},
\qquad
\Rm(u)-u=\frac{u(1-u)}{1-u+u^2}\,\bigl[\alpha u-\beta(1-u)\bigr].
\label{eq:S-onestep}
\end{equation}
Tolerance raises and strictness lowers the stationary reputation relative to
instantaneous updating.
\item Full readout, $\alpha=\beta=1$: $\Sigma_\W=(1-v^{k_1+1})/u$ and
$\Sigma_\S=0$, so
\begin{equation}
\Rm(u)=\frac{u^{k_2}\bigl(1-v^{k_1+1}\bigr)}{Z(u)},\qquad
1-\Rm(u)=\frac{v^{k_1}\bigl(1-u^{k_2+1}\bigr)}{Z(u)} .
\label{eq:S-full}
\end{equation}
Exchanging $u\leftrightarrow v$ and $k_1\leftrightarrow k_2$ exchanges
$\Rm\leftrightarrow1-\Rm$: tolerance and strictness depths act symmetrically on
the two kinds of reputational error.
\item No recovery path, $k_2=0$: $Z=1$ and
$\Rm(u)=u\sum_{m=0}^{k_1}(\alpha v)^m$; $\beta$ has no effect.
\end{enumerate}

\subsection{Interpretation through passage times}
Let $\Theta_{\G\to\B}$ be the number of assessments a record starting in $\G$
receives until it first enters $\B$, and $\Theta_{\B\to\G}$ the number from $\B$
until it first enters $\G$. Leaving good standing requires a run of $k_1+1$
negative assessments, and leaving bad standing a run of $k_2+1$ positive ones.
The classical waiting time for a run gives
\begin{equation}
\mathbb E[\Theta_{\G\to\B}]=\frac{1-v^{k_1+1}}{u\,v^{k_1+1}},\qquad
\mathbb E[\Theta_{\B\to\G}]=\frac{1-u^{k_2+1}}{v\,u^{k_2+1}} .
\label{eq:S-passage}
\end{equation}
The good-standing block is entered only at $\G$ and the bad-standing block only
at $\B$, so the record alternates between sojourns of these two kinds.
Renewal--reward then gives the stationary probability of good standing,
$\mathbb E[\Theta_{\G\to\B}]/(\mathbb E[\Theta_{\G\to\B}]+\mathbb E[\Theta_{\B\to\G}])$,
which coincides with Eq.~\eqref{eq:S-full}. With full readout, therefore,
$\Rm(u)$ is the fraction of time spent in good standing, and the depths enter
only through the two mean sojourn times. With partial readout ($\alpha<1$ or
$\beta<1$), part of each sojourn is read as the opposite image.

\subsection{Properties used in the analysis}
\begin{proposition}\label{prop:properties}
Let $0<u<1$.
\begin{enumerate}
\item[(a)] Readout monotonicity:
\begin{equation}
\partial_\alpha\Rm=\frac{u^{k_2+1}}{Z}\sum_{m=1}^{k_1}m\,\alpha^{m-1}v^m\ \ge0,
\qquad
\partial_\beta\Rm=-\frac{v^{k_1+1}}{Z}\sum_{n=1}^{k_2}n\,\beta^{n-1}u^n\ \le0 .
\label{eq:S-readout-derivatives}
\end{equation}
\item[(b)] Derivative. With $\Omega=u^{k_2+1}\Sigma_\W+v^{k_1+1}\Sigma_\S$ the numerator of Eq.~\eqref{eq:S-map},
\begin{equation}
\Rm'=\frac{\Omega'Z-\Omega Z'}{Z^2},\quad
\begin{aligned}
\Omega'&=(k_2+1)u^{k_2}\Sigma_\W-u^{k_2+1}\textstyle\sum_{m=1}^{k_1}m\alpha^mv^{m-1}
-(k_1+1)v^{k_1}\Sigma_\S+v^{k_1+1}\textstyle\sum_{n=1}^{k_2}nu^{n-1}(1-\beta^n),\\
Z'&=-k_1v^{k_1-1}(1-u^{k_2})+k_2u^{k_2-1}(1-v^{k_1}).
\end{aligned}
\label{eq:S-derivative}
\end{equation}
\item[(c)] Response of a mostly good record ($v\to0$). If $k_1\ge1$,
\begin{equation}
1-\Rm(1-v)=(1-\alpha)\,v+O(v^2),\qquad \Rm'(1)=1-\alpha .
\label{eq:S-near-one}
\end{equation}
If $\alpha=1$, for any $k_1\ge0$,
\begin{equation}
1-\Rm(1-v)=\Bigl(1+\sum_{n=1}^{k_2}\beta^n\Bigr)v^{k_1+1}+O(v^{k_1+2}),
\qquad
\Rm'(1-v)=(k_1+1)\Bigl(1+\sum_{n=1}^{k_2}\beta^n\Bigr)v^{k_1}+O(v^{k_1+1}).
\label{eq:S-near-one-full}
\end{equation}
\item[(d)] Response of a mostly bad record ($u\to0$). If $k_2\ge1$,
$\Rm(u)=(1-\beta)u+O(u^2)$. If $\beta=1$, for any $k_2\ge0$,
$\Rm(u)=\bigl(1+\sum_{m=1}^{k_1}\alpha^m\bigr)u^{k_2+1}+O(u^{k_2+2})$.
\end{enumerate}
\end{proposition}
\begin{proof}
(a) and (b) follow by differentiating Eq.~\eqref{eq:S-map}; $Z$ does not depend
on $\alpha$ or $\beta$. (c) Write
$1-\Rm=\pi_\B+\sum_m(1-\alpha^m)\pi_{\W_m}+\sum_n\beta^n\pi_{\S_n}$. As $v\to0$
with $k_1\ge1$, $Z\to1$, $\pi_{\W_1}=v\pi_\G=v+O(v^2)$, and
$\pi_\B,\pi_{\W_{m\ge2}},\pi_{\S_n}=O(v^2)$, which gives the first statement.
If $\alpha=1$ the warning terms vanish, and
$1-\Rm=\pi_\B\bigl(1+\sum_n\beta^nu^n\bigr)=v^{k_1+1}\bigl(1+\sum_n\beta^n\bigr)+O(v^{k_1+2})$;
differentiating gives the derivative. (d) follows in the same way from
$\Rm=\pi_\G\Sigma_\W+\sum_n(1-\beta^n)\pi_{\S_n}$.
\end{proof}

Parts (c) and (d) quantify the two errors that inertia filters. A good
individual that is occasionally assessed negatively (rate $v$) is read as bad
with probability $(1-\alpha)v$ as soon as a warning path exists, irrespective of
its length; only with $\alpha=1$ does this error become of order $v^{k_1+1}$.
Symmetrically, a bad individual that is occasionally assessed positively is
read as good with probability $(1-\beta)u$, or of order $u^{k_2+1}$ when $\beta=1$.
The same expansion implies that, near an all-good state, the marginal reputational
penalty for an additional negative assessment is $\Rm'(1)=1-\alpha$, which
vanishes under full tolerance readout (Supplementary Note~\ref{sec:opportunism}).

\begin{proposition}[Concavity without a recovery path]\label{prop:concave}
If $k_2=0$, $\Rm$ is concave on $[0,1]$ for every $\alpha\in[0,1]$ and $k_1\ge0$.
\end{proposition}
\begin{proof}
For $k_2=0$, $\Rm(u)=u\sum_{m=0}^{k_1}\alpha^m v^m$. Using
$u\sum_{i=0}^{j}v^i=1-v^{j+1}$,
\[
\Rm(u)=(1-\alpha)\sum_{j=0}^{k_1-1}\alpha^j\bigl[1-v^{j+1}\bigr]+\alpha^{k_1}\bigl[1-v^{k_1+1}\bigr],
\]
because the coefficient of $u\,v^i$ on the right is
$(1-\alpha)\sum_{j=i}^{k_1-1}\alpha^j+\alpha^{k_1}=\alpha^i$. The weights are
non-negative, and each $1-(1-u)^{j+1}$ is concave in $u$.
\end{proof}
Proposition~\ref{prop:concave} covers the selected SC $(2,0)$ and SH $(3,0)$
architectures. With a recovery path and $\beta=1$, $\Rm$ is not concave: it
starts as $u^{k_2+1}$ near $u=0$, and for $k_1\ge1$ and $\alpha=1$ it approaches
1 as $1-cv^{k_1+1}$ near $u=1$.

% =====================================================================
\section{Reputation dynamics in the population}
\label{sec:population}

\subsection{Reputation quantities}
At fixed strategy composition $\boldsymbol f$, let $r_I$ be the probability that
a randomly chosen observer regards a randomly chosen type-$I$ individual as
good, that is, the expected readout probability $\sigma(L_{ij})$ over observers
$i$ and type-$I$ targets $j$. The probability that a random recipient is
regarded as good is $\bar r=\sum_If_Ir_I$. Because all individuals use the same
norm and record mechanism, all observers are statistically equivalent in the
mean-field description; in the finite simulations, only DISC observers are
stored because only they act on reputations (Supplementary
Note~\ref{sec:simulations}).

\subsection{Mean-field closure}
Two assumptions turn the single-record result of Supplementary
Note~\ref{sec:map} into population equations; both are standard in analyses
of private assessment~\cite{okada2018,radzvilavicius2019,kessinger2023,michelmata2024,kawakatsu2024}.
\begin{enumerate}
\item[(A1)] \emph{Conditional independence of images.} The images that the donor
and the observer hold of a recipient are independent given the recipient's
strategy.
\item[(A2)] \emph{Independent inputs.} The successive assessments entering any
record about a type-$I$ target are independent, each positive with the same
probability $u_I$.
\end{enumerate}
(A1) neglects correlations between the records that different observers hold
of the same individual. Such correlations arise because two observers may have
assessed the same past actions of that individual, and because targets of the
same strategy differ in their histories. (A2) neglects temporal correlations
in the inputs to a record. The finite-population simulations retain both
effects and are used to test reputation formation without the closure.

\subsection{Positive-assessment probabilities}
Under (A1), donor and observer both regard a random recipient as good, both as
bad, or hold a given ordered disagreement with probabilities
\begin{equation}
h_{\G\G}=\sum_Jf_Jr_J^2,\qquad
h_{\B\B}=\sum_Jf_J(1-r_J)^2,\qquad
h_{\G\B}=\sum_Jf_Jr_J(1-r_J),
\label{eq:S-pairs}
\end{equation}
with $h_{\G\G}+2h_{\G\B}+h_{\B\B}=1$. Conditioning on the recipient's strategy
matters: $h_{\G\G}\ge\bar r^{\,2}$, with equality only if all strategies have the
same reputation. An ALLC donor intends $\C$ whatever the recipient, and an ALLD
donor intends $\D$; their assessments depend only on the observer's image.
A DISC donor intends $\C$ if its own image is $\G$; Supplementary
Table~\ref{tab:disc-cases} lists the four cases. Summing,
\begin{equation}
u_{\ALLC}=\Phi_\C(\bar r),\qquad u_{\ALLD}=\Phi_\D(\bar r),\qquad
u_{\DISC}=h_{\G\G}\Pp CG+h_{\G\B}\bigl(\Pp CB+\Pp DG\bigr)+h_{\B\B}\Pp DB .
\label{eq:S-u}
\end{equation}
Using Eq.~\eqref{eq:S-uR}, the DISC expression can be written as
$u_{\DISC}=\sum_Jf_J\,u_{\mathrm R}(r_J)$: a DISC donor is assessed as if it were
a resident in a homogeneous population whose reputation equals that of its
recipient, averaged over recipients.

\begin{table}[htbp]
\centering\small
\caption{\textbf{Assessment of a DISC donor.} The first image is the donor's, the second the observer's.}
\label{tab:disc-cases}
\begin{tabular}{@{}llll@{}}
\toprule
Images (donor, observer) & Probability & Intended action & Positive assessment\\
\midrule
$(\G,\G)$ & $h_{\G\G}$ & $\C$ & $\Pp CG$\\
$(\G,\B)$ & $h_{\G\B}$ & $\C$ & $\Pp CB$\\
$(\B,\G)$ & $h_{\G\B}$ & $\D$ & $\Pp DG$\\
$(\B,\B)$ & $h_{\B\B}$ & $\D$ & $\Pp DB$\\
\bottomrule
\end{tabular}
\end{table}

\subsection{Self-consistency}
Under (A2), each record about a type-$I$ target has the stationary response
$\Rm(u_I)$. Stationary reputations therefore satisfy
\begin{equation}
r_I^*=\Rm\bigl(u_I(\boldsymbol r^*;\boldsymbol f)\bigr),\qquad
I\in\{\ALLC,\ALLD,\DISC\}.
\label{eq:S-self-consistency}
\end{equation}
Reputations determine assessments through Eq.~\eqref{eq:S-u}, and assessments
determine reputations through $\Rm$. Because $u_{\ALLC}$ and $u_{\ALLD}$ depend on
$\boldsymbol r$ only through $\bar r$, Eq.~\eqref{eq:S-self-consistency} is a
three-dimensional fixed-point problem that can have several solutions.

\subsection{Full latent-state dynamics and the stationary reduction}
Within the same closure, let $\boldsymbol\pi_I(t')$ be the row vector of
latent-state probabilities of records about type-$I$ targets, with time $t'$
measured in assessments per record. Asynchronous updating, in which each record
is updated at unit rate, gives
\begin{equation}
\frac{\mathrm d\boldsymbol\pi_I}{\mathrm dt'}
=\boldsymbol\pi_I\bigl\{T\bigl[u_I(\boldsymbol r;\boldsymbol f)\bigr]-\mathbf I\bigr\},
\qquad r_I=\boldsymbol\pi_I\boldsymbol\sigma .
\label{eq:S-full-state}
\end{equation}
This flow conserves probability and non-negativity. It is the mean-field limit
of the finite simulation when time is measured in sweeps (one update per stored
record on average) and assessment probabilities are replaced by their closed
forms.

\begin{proposition}[Stationary equivalence]\label{prop:equivalence}
Suppose $0<u_I<1$ at a stationary state. Then $\{\boldsymbol\pi_I^*\}$ is a
stationary state of Eq.~\eqref{eq:S-full-state} if and only if
$\boldsymbol\pi_I^*=\boldsymbol\pi(u_I^*)$ with $u_I^*=u_I(\boldsymbol r^*;\boldsymbol f)$
and $\boldsymbol r^*$ solves Eq.~\eqref{eq:S-self-consistency}.
\end{proposition}
\begin{proof}
At a stationary state, $\boldsymbol\pi_I^*T(u_I^*)=\boldsymbol\pi_I^*$, so
$\boldsymbol\pi_I^*$ is the unique stationary distribution of $T(u_I^*)$
(Proposition~\ref{prop:stationary}), and
$r_I^*=\boldsymbol\pi_I^*\boldsymbol\sigma=\Rm(u_I^*)$. Conversely, given a solution
$\boldsymbol r^*$ of Eq.~\eqref{eq:S-self-consistency}, the distributions
$\boldsymbol\pi(u_I^*)$ are stationary and reproduce $\boldsymbol r^*$.
\end{proof}

Transient dynamics are not reduced in the same way. Projecting
Eq.~\eqref{eq:S-full-state} gives
$\mathrm dr_I/\mathrm dt'=\boldsymbol\pi_IT(u_I)\boldsymbol\sigma-r_I$, and
$\boldsymbol\pi_IT(u_I)\boldsymbol\sigma$ (the expected readout after one more
update) generally differs from $\Rm(u_I)$. For example, for one-step inertia
with $\alpha=\beta=1$ and all records in $\G$,
$\boldsymbol\pi_IT(u)\boldsymbol\sigma=1$ (both $\G$ and $\W_1$ are read as good),
whereas $\Rm(u)=u^2(2-u)/(1-u+u^2)<1$. Two populations with the same $r_I$ but
different amounts of accumulated evidence evolve differently.

\subsection{Auxiliary relaxation and branch selection}
To select a stationary solution we use the auxiliary flow
\begin{equation}
\frac{\mathrm dr_I}{\mathrm d\tau}=\Rm\bigl(u_I(\boldsymbol r;\boldsymbol f)\bigr)-r_I,
\label{eq:S-relaxation}
\end{equation}
whose fixed points are exactly the solutions of
Eq.~\eqref{eq:S-self-consistency}. Its Jacobian is
$\mathrm{diag}[\Rm'(u_I)]\,\partial\boldsymbol u/\partial\boldsymbol r-\mathbf I$. The
time $\tau$ is a computational device; Eq.~\eqref{eq:S-relaxation} is not the
transient dynamics of the records, and its stability need not coincide with
that of Eq.~\eqref{eq:S-full-state} in general. Three selection protocols are
used.
\begin{enumerate}
\item \emph{Main convention} (all strategy-dynamics results): at every
composition, Eq.~\eqref{eq:S-relaxation} is integrated from $r_I(0)=1$, which
corresponds to all records in $\G$. The replicator field is thus defined on the
branch reached from good standing.
\item \emph{Highest strictly stable root} (homogeneous DISC: Extended Data
Figs.~\EDResident--\EDPlanes, local invasion and Fig.~\FigOpportunism): the
highest root of $F(r)$ with $F'(r)<0$. It coincides with protocol 1 unless the
highest root is a fold, which is marginal.
\item \emph{Quasi-static tracking} (Extended Data Fig.~\EDResident e):
reputations are initialized once, at $t=0$, as all-good or all-bad and then
follow the strategy trajectory, Eq.~\eqref{eq:S-relaxation} being re-solved from
the current reputations at every evaluation. A branch is then left only where it
ends at a fold, so the outcome can depend on reputational history.
\end{enumerate}

\subsection{Homogeneous discriminator populations}
\label{sec:homogeneous}
In a homogeneous DISC population, $h_{\G\G}=r^2$, $h_{\G\B}=r(1-r)$ and
$h_{\B\B}=(1-r)^2$, and Eq.~\eqref{eq:S-u} reduces to Eq.~\eqref{eq:S-uR}. The
stationary resident reputations are the roots of
\begin{equation}
F(r)=\Rm\bigl[u_{\mathrm R}(r)\bigr]-r,\qquad
F'(r)=\Rm'\bigl[u_{\mathrm R}(r)\bigr]\,u_{\mathrm R}'(r)-1 .
\label{eq:S-F}
\end{equation}
A simple root is stable in the scalar flow $\mathrm dr/\mathrm d\tau=F(r)$ if
$F'<0$ and unstable if $F'>0$. Because $F(0)\ge0$ and $F(1)\le0$, roots
alternate in stability, and a pair of roots appears or disappears at a fold,
$F=F'=0$. Bistability (Extended Data Figs.~\EDResident{} and~\EDPlanes) means
two stable roots separated by an unstable one; with errors, $F(1)<0$, so the
highest root is approached from $r=1$. For every root in the full-readout depth
planes ($0\le k_1,k_2\le8$), the classification by $F'$ agreed with the leading
eigenvalue of Eq.~\eqref{eq:S-full-state} linearized about
$\boldsymbol\pi(u_{\mathrm R}(r))$ on the probability simplex; this linearization
includes the dependence of $u_{\mathrm R}$ on the current readout.

\emph{Sensitivity to readout.} On a regular stable branch ($F'<0$), implicit
differentiation of $F(r^*;\theta)=0$ for $\theta\in\{\alpha,\beta\}$ gives
\begin{equation}
\frac{\partial r^*}{\partial\theta}=-\frac{\partial_\theta\Rm\bigl[u_{\mathrm R}(r^*)\bigr]}{F'(r^*)},
\label{eq:S-branch-sensitivity}
\end{equation}
so by Eq.~\eqref{eq:S-readout-derivatives} the resident reputation increases
with $\alpha$ and decreases with $\beta$ along a stable branch. The derivative
diverges at a fold, where the branch ends.

\subsection{The ALLD--DISC edge}
\label{sec:edge}
On the edge $f_{\ALLC}=0$, $f_{\DISC}=x$, $f_{\ALLD}=1-x$, the three-dimensional
problem reduces exactly to one equation. Given $\bar r$,
$r_{\ALLD}=\Rm[\Phi_\D(\bar r)]$ and $r_{\ALLC}=\Rm[\Phi_\C(\bar r)]$ follow
directly, and $r_{\DISC}=[\bar r-(1-x)r_{\ALLD}]/x$ follows from the definition of
$\bar r$. Substituting these into
$r_{\DISC}=\Rm\bigl[x\,u_{\mathrm R}(r_{\DISC})+(1-x)\,u_{\mathrm R}(r_{\ALLD})\bigr]$,
which is $r_{\DISC}=\Rm(u_{\DISC})$ written with $u_{\DISC}=\sum_Jf_J\,u_{\mathrm R}(r_J)$,
gives one equation in $\bar r$,
whose roots with $r_{\DISC}\in[0,1]$ are all self-consistent states on the edge.
Their stability is that of Eq.~\eqref{eq:S-relaxation}. The normalized payoff
advantage of DISC over ALLD on the edge is
$(\Pi_{\DISC}-\Pi_{\ALLD})/[c(1-e_1)]=(b/c)\,x\,(r^*_{\DISC}-r^*_{\ALLD})-\bar r^*$
(Supplementary Note~\ref{sec:strategy}). Where it changes sign continuously, the
edge has an equilibrium; where the selected branch ends at a fold, the sign can
change discontinuously (Extended Data Fig.~\EDEdge).

% =====================================================================
\section{Strategy dynamics and cooperation}
\label{sec:strategy}

\subsection{Payoffs}
Reputations are assumed to reach their stationary state quickly compared with
changes in strategy frequencies, as in previous models of indirect
reciprocity~\cite{radzvilavicius2019,kessinger2023,michelmata2024,radzvilavicius2021}.
Payoffs are computed per unit time, in which each individual acts on average
once as donor and once as recipient. A type-$I$ recipient receives intended
help from ALLC donors (probability $f_{\ALLC}$) and from DISC donors that regard
it as good (probability $f_{\DISC}r_I^*$). As a donor, ALLC intends to help with
probability 1, ALLD with probability 0, and DISC with probability $\bar r^*$,
because its recipient is a random individual. Each intended donation is
realized with probability $1-e_1$, so
\begin{equation}
\Pi_I=(1-e_1)\Bigl[b\,\bigl(f_{\ALLC}+f_{\DISC}r_I^*\bigr)-c\,\chi_I\Bigr],\qquad
\chi_{\ALLC}=1,\ \chi_{\ALLD}=0,\ \chi_{\DISC}=\bar r^*,
\label{eq:S-payoffs}
\end{equation}
which gives the payoffs of the Methods. Every realized donation creates one
benefit and one cost, so $\bar\Pi=\sum_If_I\Pi_I=(b-c)\rho$, where
$\rho=(1-e_1)\rho_{\mathrm{int}}$ and $\rho_{\mathrm{int}}=f_{\ALLC}+f_{\DISC}\bar r^*$.
The factor $1-e_1$ rescales time in Eq.~\eqref{eq:S-replicator}, but $e_1$ also
changes the reputations through $P^+_{A\mid Y}$.

\subsection{Replicator dynamics and equilibria}
Strategy frequencies follow
\begin{equation}
\frac{\mathrm df_I}{\mathrm dt}=f_I\bigl(\Pi_I-\bar\Pi\bigr),
\label{eq:S-replicator}
\end{equation}
with stationary reputations recomputed at every composition. Edges and
vertices are invariant. At a vertex where all individuals use $J$, the
eigenvalues are the growth rates $\Pi_I-\Pi_J$ of the two absent strategies,
evaluated with the reputations of that vertex; at the DISC vertex these are the
invasion payoffs of Supplementary Note~\ref{sec:invasion}. At an equilibrium on
an edge, the tangent eigenvalue describes stability within the edge and the
transverse eigenvalue is the growth rate of the absent strategy. Interior
equilibria are classified by the Jacobian of Eq.~\eqref{eq:S-replicator} on the
simplex. Because reputations can jump where the selected branch ends, the
replicator field can be discontinuous at folds (Fig.~\FigStrategy{} and
Extended Data Fig.~\EDEdge).

\begin{proposition}[Neutral line for Scoring without inertia]\label{prop:sc-line}
Under Scoring with $\Rm(u)=u$, $\Pi_{\ALLC}=\Pi_{\ALLD}=\Pi_{\DISC}$ at every
composition with $f_{\DISC}=c/[b(\epsilon-e_2)]$.
\end{proposition}
\begin{proof}
Under SC, $\Phi_\C\equiv\epsilon$ and $\Phi_\D\equiv e_2$ do not depend on the
recipient, so $u_{\DISC}=\epsilon(h_{\G\G}+h_{\G\B})+e_2(h_{\G\B}+h_{\B\B})=e_2+(\epsilon-e_2)\bar r$,
and $r_{\ALLC}=\epsilon$, $r_{\ALLD}=e_2$, $r_{\DISC}=e_2+(\epsilon-e_2)\bar r$. Then
$\Pi_{\ALLC}-\Pi_{\DISC}=(1-e_1)(1-\bar r)[b f_{\DISC}(\epsilon-e_2)-c]$ and
$\Pi_{\ALLD}-\Pi_{\DISC}=(1-e_1)\bar r[c-bf_{\DISC}(\epsilon-e_2)]$, which both
vanish at the stated $f_{\DISC}$.
\end{proof}
At the reference parameters, $f_{\DISC}=1/(5\times0.9408)\approx0.2126$. This line
of equilibria, neutral along its direction, is the dashed line in
Fig.~\FigStrategy a (top). Trajectories that reach it drift neutrally along it,
which is why integration under SC without inertia sometimes reaches the time
cap without affecting the outcome class.

\subsection{Cooperation measures}
At a composition, intended and realized cooperation are $\rho_{\mathrm{int}}$ and
$\rho$; in a homogeneous DISC population, $\rho_{\mathrm{int}}=r^*_{\mathrm R}$. To
compare settings globally, trajectories are integrated from $n_{\mathrm{init}}$
area-uniform initial compositions. Across all scans, every final state belonged
to one of two classes: the zero-cooperation set, comprising the ALLD vertex
together with the low-reputation segment of the ALLD--DISC edge, on which DISC
regard one another as bad, earn slightly less than ALLD and drift slowly
towards ALLD (final $\rho\le7\times10^{-4}$); or an attractor with positive
cooperation (final $\rho\ge0.34$; between 0.974 and 0.980 in nearly all
conditions). The cooperative basin is the set of initial compositions whose
trajectories reach the second class. A trajectory is classified as cooperative
if its final $\rho\ge0.01$; every threshold in the observed gap gives the same
classification, and the gap was checked for every condition. With
$\bar\rho_{\D}$ the mean final cooperation of the remaining trajectories,
\begin{equation}
\bar\rho=\omega_\C\,\bar\rho_\C+(1-\omega_\C)\,\bar\rho_{\D}\approx\omega_\C\,\bar\rho_\C ,
\label{eq:S-basin}
\end{equation}
so a change in $\bar\rho$ can be attributed to the basin size and to the
cooperation level within it. At the reference parameters, $\bar\rho_\C$ is
0.975--0.980 for one-step and selected multi-step inertia under all norms,
limited by implementation error, and at least 94\% of the increase in $\bar\rho$
produced by the selected architectures comes from the enlarged basin. These
quantities depend on the area-uniform sampling of initial compositions; they
are not probabilities over any other distribution of starting points.

% =====================================================================
\section{Invasion of a discriminator population by unconditional strategies}
\label{sec:invasion}

\subsection{A rare mutant in a discriminator population}
Consider a homogeneous DISC population whose reputation is the selected root
$\rR$, and a rare mutant X that uses the residents' norm and record mechanism.
Because X is rare, its recipients, its observers and the individuals it is
assessed against are residents. Residents regard each other as good with
probability $\rR$, and their assessment environment is unchanged.

\begin{lemma}[Rare-mutant identity]\label{lem:identity}
Suppose that X intends to help a resident recipient with probability
$\chi_{\mathrm X}$, independently of the observer's image of that recipient and
independently across interactions. Then X is assessed positively with
probability
\begin{equation}
u_{\mathrm X}=\chi_{\mathrm X}\Phi^*_\C+(1-\chi_{\mathrm X})\Phi^*_\D
=\uR-(\rR-\chi_{\mathrm X})\,\Delta\Phi^*,
\label{eq:S-uX}
\end{equation}
its stationary reputation is $r^*_{\mathrm X}=\Rm(u_{\mathrm X})$, and
\begin{equation}
\Lambda_{\mathrm X}\equiv\frac{\Pi_{\mathrm X}-\Pi_{\mathrm R}}{c(1-e_1)}
=\bigl(\rR-\chi_{\mathrm X}\bigr)-\frac bc\bigl(\rR-r^*_{\mathrm X}\bigr).
\label{eq:S-identity}
\end{equation}
\end{lemma}
\begin{proof}
The observer regards the resident recipient as good with probability $\rR$,
independently of X's intention, so X's assessment probability is the mixture of
$\Phi_\C^*$ and $\Phi_\D^*$ with weights $\chi_{\mathrm X}$ and $1-\chi_{\mathrm X}$.
Independent actions imply independent assessments, so (A2) holds exactly for
X's records in the fixed resident environment, and
$r^*_{\mathrm X}=\Rm(u_{\mathrm X})$. A resident receives help from residents with
probability $\rR$ and helps with probability $\rR$:
$\Pi_{\mathrm R}=(1-e_1)(b-c)\rR$. The mutant receives help with probability
$r^*_{\mathrm X}$ and helps with probability $\chi_{\mathrm X}$:
$\Pi_{\mathrm X}=(1-e_1)(br^*_{\mathrm X}-c\chi_{\mathrm X})$. Subtracting and dividing
by $c(1-e_1)>0$ gives Eq.~\eqref{eq:S-identity}. The second form of
Eq.~\eqref{eq:S-uX} uses $\uR=\rR\Phi^*_\C+(1-\rR)\Phi^*_\D$.
\end{proof}

The first term of Eq.~\eqref{eq:S-identity} is the cooperation cost the mutant
saves relative to a resident; the second is the help it loses because it is
regarded differently. The resident corresponds to $\chi=\rR$, ALLD to
$\chi=0$, ALLC to $\chi=1$ and the opportunistic discriminator of
Supplementary Note~\ref{sec:opportunism} to $\chi=(1-\delta)\rR$. All of these
mutants therefore lie on the single curve
$\chi\mapsto\Rm[\uR-(\rR-\chi)\Delta\Phi^*]$, which is $\Rm$ evaluated along the
segment between $\Phi^*_\D$ ($\chi=0$) and $\Phi^*_\C$ ($\chi=1$), passing
through $\uR$ at $\chi=\rR$.

\subsection{ALLD and ALLC}
Write $r^*_{\ALLD}=\Rm(\Phi^*_\D)$ and $r^*_{\ALLC}=\Rm(\Phi^*_\C)$, with
\begin{equation}
\Phi^*_\D=\rR\Pp DG+(1-\rR)\Pp DB,\qquad
\Phi^*_\C=\rR\Pp CG+(1-\rR)\Pp CB .
\label{eq:S-PhiStar}
\end{equation}
These are the reputations of rare mutants in the resident population, not of
homogeneous ALLC or ALLD populations; they equal $r^*_{\ALLC}$ and $r^*_{\ALLD}$
of the population equations evaluated at the DISC vertex. Lemma~\ref{lem:identity}
gives
\begin{equation}
\lambda_{\ALLD}=(1-e_1)\bigl[c\,\rR-b\,(\rR-r^*_{\ALLD})\bigr],\qquad
\lambda_{\ALLC}=(1-e_1)\bigl[b\,(r^*_{\ALLC}-\rR)-c\,(1-\rR)\bigr],
\label{eq:S-lambda}
\end{equation}
where $\lambda_X=\Pi_X-\Pi_{\mathrm R}$. ALLD saves the resident's donation cost
$c\rR$ but may be helped less; ALLC can be regarded better than residents because
it also helps recipients that residents regard as bad, but pays for these
donations.

\begin{proposition}[Resistance to ALLC and ALLD]\label{prop:window}
\begin{enumerate}
\item[(a)] If $\rR>r^*_{\ALLD}$, rare ALLD is excluded ($\lambda_{\ALLD}<0$) if and
only if $b/c>R_{\mathrm L}=\rR/(\rR-r^*_{\ALLD})$. If $\rR\le r^*_{\ALLD}$ and
$\rR>0$, no $b/c$ excludes ALLD.
\item[(b)] If $r^*_{\ALLC}>\rR$, rare ALLC is excluded ($\lambda_{\ALLC}<0$) if and
only if $b/c<R_{\mathrm U}=(1-\rR)/(r^*_{\ALLC}-\rR)$. If $r^*_{\ALLC}\le\rR<1$,
ALLC is excluded for every $b/c$.
\item[(c)] If $\rR>r^*_{\ALLD}$ and $r^*_{\ALLC}>\rR$, the residents resist both
mutants if and only if
\begin{equation}
\frac{\rR}{\rR-\Rm\bigl[\rR\Pp DG+(1-\rR)\Pp DB\bigr]}
<\frac bc<
\frac{1-\rR}{\Rm\bigl[\rR\Pp CG+(1-\rR)\Pp CB\bigr]-\rR}.
\label{eq:S-window}
\end{equation}
\end{enumerate}
\end{proposition}
\begin{proof}
Rearrange the signs of Eq.~\eqref{eq:S-lambda}: $\lambda_{\ALLD}<0$ if and only if
$b(\rR-r^*_{\ALLD})>c\rR$, and $\lambda_{\ALLC}<0$ if and only if
$b(r^*_{\ALLC}-\rR)<c(1-\rR)$. Divide by the positive quantities stated in each
case.
\end{proof}

\subsection{When is the window open?}
\begin{proposition}[Chord condition]\label{prop:chord}
Under the conditions of Proposition~\ref{prop:window}(c), the interval in
Eq.~\eqref{eq:S-window} is non-empty if and only if
\begin{equation}
\rR\,\Rm(\Phi^*_\C)+(1-\rR)\,\Rm(\Phi^*_\D)<\Rm(\uR),
\qquad \uR=\rR\Phi^*_\C+(1-\rR)\Phi^*_\D .
\label{eq:S-chord}
\end{equation}
\end{proposition}
\begin{proof}
$R_{\mathrm L}<R_{\mathrm U}$ is equivalent to
$\rR(r^*_{\ALLC}-\rR)<(1-\rR)(\rR-r^*_{\ALLD})$, that is,
$\rR r^*_{\ALLC}+(1-\rR)r^*_{\ALLD}<\rR$. Self-consistency gives
$\rR=\Rm(\uR)$.
\end{proof}
The resident's own action is a mixture of the actions of ALLC and ALLD with
weights $\rR$ and $1-\rR$, and its assessment probability $\uR$ is the
corresponding mixture of $\Phi^*_\C$ and $\Phi^*_\D$. The window is open
precisely when the reputation of this mixture exceeds the same mixture of the
two pure reputations, that is, when $\Rm$ lies above its chord between
$\Phi^*_\D$ and $\Phi^*_\C$ at $\uR$.

\begin{corollary}[No window without inertia]\label{cor:nowindow}
If $\Rm(u)=u$ and $\Delta\Phi^*>0$, then $R_{\mathrm L}=R_{\mathrm U}=1/\Delta\Phi^*$:
resistance to ALLD requires $b/c>1/\Delta\Phi^*$ and resistance to ALLC requires
$b/c<1/\Delta\Phi^*$, so no benefit-to-cost ratio gives resistance to both.
\end{corollary}
\begin{proof}
With $\Rm(u)=u$, $\rR=\uR$, $\rR-r^*_{\ALLD}=\rR\Delta\Phi^*$ and
$r^*_{\ALLC}-\rR=(1-\rR)\Delta\Phi^*$.
\end{proof}
If $\Delta\Phi^*\le0$ without inertia (SJ, where the instantaneous resident has
$\rR=1/2$ and $\Delta\Phi^*=0$), ALLD is regarded at least as well as the
residents and cannot be excluded. The thresholds also have a marginal
interpretation. Define the reputation lost per unit of withheld cooperation,
$H_{\ALLD}=(\rR-r^*_{\ALLD})/\rR$, and the reputation gained per unit of extra
cooperation, $H_{\ALLC}=(r^*_{\ALLC}-\rR)/(1-\rR)$. Then
$R_{\mathrm L}=1/H_{\ALLD}$, $R_{\mathrm U}=1/H_{\ALLC}$, and the window is
\begin{equation}
H_{\ALLC}<\frac cb<H_{\ALLD}:
\label{eq:S-marginal-window}
\end{equation}
cooperating more than the residents must not pay off, and cooperating less
must not pay off. Because $\Phi^*_\C-\uR=(1-\rR)\Delta\Phi^*$ and
$\uR-\Phi^*_\D=\rR\Delta\Phi^*$, $H_{\ALLC}$ and $H_{\ALLD}$ are $\Delta\Phi^*$ times
the slopes of the chords of $\Rm$ on $[\uR,\Phi^*_\C]$ and $[\Phi^*_\D,\uR]$,
respectively. Inertia opens the window by making $\Rm$ flatter above the
resident's assessment probability than below it, so that extra cooperation buys
little additional reputation while withheld cooperation is still penalized.
Supplementary Table~\ref{tab:reference} lists the values at the reference
parameters; Supplementary Note~\ref{sec:numerics} gives the numerical
classification rule.

\begin{table}[htbp]
\centering\small
\caption{\textbf{Reference values for a resident discriminator population.} Highest stable homogeneous-DISC root $\rR$, reputations of rare ALLD and ALLC, and the thresholds of Eq.~\eqref{eq:S-window}; $e_1=e_2=0.02$, $\alpha=\beta=1$ for inertial settings. Without inertia $R_{\mathrm L}=R_{\mathrm U}$ (Corollary~\ref{cor:nowindow}); under SJ without inertia ALLD is regarded as well as residents and cannot be excluded. Values agree with the verified data to all digits shown.}
\label{tab:reference}
\begin{tabular}{@{}llccccc@{}}
\toprule
Norm & Setting & $\rR$ & $r^*_{\ALLD}$ & $r^*_{\ALLC}$ & $R_{\mathrm L}$ & $R_{\mathrm U}$\\
\midrule
SC & none $(0,0)$ & 0.337838 & 0.020000 & 0.960800 & 1.0629 & 1.0629\\
   & one-step $(1,1)$ & 0.996279 & 0.000808 & 0.996869 & 1.0008 & 6.3070\\
   & selected $(2,0)$ & 0.999940 & 0.058808 & 0.999940 & 1.0625 & 230.24\\
SH & none $(0,0)$ & 0.020391 & 0.020000 & 0.039184 & 52.127 & 52.127\\
   & one-step $(1,1)$ & 0.995272 & 0.000808 & 0.996111 & 1.0008 & 5.6409\\
   & selected $(3,0)$ & 0.999998 & 0.077632 & 0.999998 & 1.0842 & 4410.4\\
SJ & none $(0,0)$ & 0.500000 & 0.500000 & 0.500000 & -- & --\\
   & one-step $(1,1)$ & 0.995307 & 0.001215 & 0.996133 & 1.0012 & 5.6865\\
   & selected $(4,3)$ & 0.9999996 & $8.3\times10^{-7}$ & 0.9999996 & 1.0000 & 20910\\
SS & none $(0,0)$ & 0.833333 & 0.180000 & 0.964000 & 1.2755 & 1.2755\\
   & one-step $(1,1)$ & 0.996297 & 0.001122 & 0.996880 & 1.0011 & 6.3441\\
   & selected $(1,2)$ & 0.993660 & 0.000036 & 0.995245 & 1.0000 & 4.0001\\
\bottomrule
\end{tabular}
\end{table}

% =====================================================================
\section{Opportunistic invasion and secure readout}
\label{sec:opportunism}

\subsection{The opportunistic discriminator}
An opportunistic discriminator (O-DISC) uses the residents' norm and record
mechanism and follows DISC, except that whenever its own image of the
recipient is $\G$ it independently replaces the intended cooperation by
intended defection with probability $\delta\in[0,1]$. This choice precedes
implementation error, and O-DISC never cooperates with recipients it regards as
bad. Its images of resident recipients are good with probability $\rR$ and are
independent of the observer's images by (A1), so it intends to help with
probability $\chi=(1-\delta)\rR$. By Lemma~\ref{lem:identity},
\begin{equation}
u_{\mathrm X}(\delta)=(1-\delta)\rR\,\Phi^*_\C+\bigl[1-(1-\delta)\rR\bigr]\Phi^*_\D
=\uR-\delta\,\rR\,\Delta\Phi^*,
\qquad
r^*_{\mathrm X}(\delta)=\Rm\bigl[u_{\mathrm X}(\delta)\bigr],
\label{eq:S-odisc-u}
\end{equation}
\begin{equation}
\Lambda(\delta)=\frac{\Pi_{\mathrm X}-\Pi_{\mathrm R}}{c(1-e_1)}
=\delta\,\rR-\frac bc\bigl[\rR-r^*_{\mathrm X}(\delta)\bigr].
\label{eq:S-Lambda}
\end{equation}
Equivalently, $u_{\mathrm X}(\delta)=(1-\delta)\uR+\delta\Phi^*_\D$: the mutant's
assessment probability moves linearly from the resident's ($\delta=0$) to that
of ALLD ($\delta=1$). Hence $\Lambda(0)=0$, the mutant being a neutral copy of
a resident, and $\Lambda(1)=\lambda_{\ALLD}/[c(1-e_1)]$. The family contains
neither ALLC nor mutants whose omissions are correlated in time.

\subsection{Weak opportunism}
Differentiating Eq.~\eqref{eq:S-Lambda}, with
$\mathrm dr^*_{\mathrm X}/\mathrm d\delta=-\rR\Delta\Phi^*\Rm'[u_{\mathrm X}(\delta)]$,
\begin{equation}
\Lambda'(0)=\rR\Bigl[1-\frac bc\,\Delta\Phi^*\,\Rm'(\uR)\Bigr].
\label{eq:S-slope}
\end{equation}
If $\Lambda'(0)>0$, all sufficiently weak opportunists invade; if
$\Lambda'(0)<0$, they are excluded, but larger $\delta$ may still invade. The
product $\Delta\Phi^*\Rm'(\uR)$ is the reputation lost per unit of withheld
cooperation for an infinitesimal deviation: the norm converts withheld
cooperation into a lower positive-assessment probability at rate
$\Delta\Phi^*$, and the record converts this into a lower reputation at rate
$\Rm'(\uR)$.

\subsection{Resistance to the whole family}
For $\rR>0$ and $0<\delta\le1$, define the reputation lost per unit of withheld
cooperation,
\begin{equation}
H(\delta)=\frac{\rR-r^*_{\mathrm X}(\delta)}{\delta\,\rR}
=\Delta\Phi^*\,\frac{\Rm(\uR)-\Rm[u_{\mathrm X}(\delta)]}{\uR-u_{\mathrm X}(\delta)}
=\Delta\Phi^*\int_0^1\Rm'\bigl(\uR-s\,\delta\,\rR\,\Delta\Phi^*\bigr)\,\mathrm ds ,
\label{eq:S-H}
\end{equation}
and its continuous extension $H(0)=\Delta\Phi^*\Rm'(\uR)$. The second form uses
$\rR=\Rm(\uR)$ and $\uR-u_{\mathrm X}=\delta\rR\Delta\Phi^*$ (for $\Delta\Phi^*\ne0$):
$H(\delta)$ is $\Delta\Phi^*$ times the slope of the chord of $\Rm$ between
$u_{\mathrm X}(\delta)$ and $\uR$. Then
$\Lambda(\delta)=\delta\rR[1-(b/c)H(\delta)]$, and $H(1)=H_{\ALLD}$.

\begin{proposition}[Family criterion]\label{prop:family}
Let $\rR>0$ and $H_{\min}=\min_{0\le\delta\le1}H(\delta)$. Then
$\Lambda(\delta)\le0$ for every $\delta\in(0,1]$ if and only if
$H_{\min}\ge c/b$. If $H_{\min}>c/b$, every O-DISC mutant is strictly excluded;
if $H_{\min}<c/b$, a mutant with $\delta$ in a non-empty interval invades.
\end{proposition}
\begin{proof}
For $\delta>0$, $\Lambda(\delta)$ has the sign of $1-(b/c)H(\delta)$. The
minimum includes $\delta=0$ through the continuous extension, which is required
because $\Lambda'(0)$ has the sign of $1-(b/c)H(0)$; if $H(0)<c/b$, all small
$\delta>0$ invade although no single $\delta>0$ attains $H(0)$. Continuity of
$H$ gives the interval.
\end{proof}

\begin{corollary}[Curvature and the location of the worst mutant]\label{cor:curvature}
Let $\Delta\Phi^*>0$.
\begin{enumerate}
\item[(a)] If $\Rm$ is concave on $[\Phi^*_\D,\uR]$, then $H$ is non-decreasing in
$\delta$, $H_{\min}=H(0)$, and the family is excluded if and only if
$\Lambda'(0)\le0$. This holds for every architecture with $k_2=0$
(Proposition~\ref{prop:concave}).
\item[(b)] If $\Rm$ is convex on $[\Phi^*_\D,\uR]$, then $H_{\min}=H(1)$ and
resistance to ALLD implies resistance to the whole family.
\item[(c)] In general, the worst mutant can be interior, and all $\delta$ must be
checked.
\end{enumerate}
\end{corollary}
\begin{proof}
As $\delta$ increases, $u_{\mathrm X}(\delta)$ decreases from $\uR$ to $\Phi^*_\D$.
For concave $\Rm$, the slope of the chord from $u_{\mathrm X}$ to the fixed right
endpoint $\uR$ is non-increasing in $u_{\mathrm X}$, hence non-decreasing in
$\delta$; for convex $\Rm$ it is non-increasing in $\delta$.
\end{proof}
Part (a) also shows how tolerance creates vulnerability. With a concave
response, the reputation lost by a mutant is at most the linear interpolation
between the resident and ALLD, $\rR-r^*_{\mathrm X}(\delta)\le\delta(\rR-r^*_{\ALLD})$,
so $\Lambda(\delta)\ge\delta\,\Lambda(1)$. Excluding ALLD then says little about
partial opportunism. Near an all-good resident, Eq.~\eqref{eq:S-near-one} gives
$\Rm'(\uR)\approx1-\alpha$ for $k_1\ge1$; with full tolerance readout,
Eq.~\eqref{eq:S-near-one-full} makes $\Rm'(\uR)$ of order $(1-\uR)^{k_1}$, so
$H(0)$ and the marginal penalty vanish as the warning path lengthens. For SJ
$(4,3)$ with $\beta=1$ at the reference parameters, $H(0)=4.8\times10^{-5}$ at
$\alpha=1$ but $0.2037$ at $\alpha=0.8$, where $c/b=0.2$.

\subsection{An exact threshold in the error-free limit}
\begin{proposition}[Error-free security threshold]\label{prop:errorfree}
Let $e_1=e_2=0$, $\beta=1$ and $k_1\ge1$, and consider the all-good resident
$\rR=1$ under any of the four norms and any $k_2\ge0$. Every O-DISC mutant is
excluded, $\Lambda(\delta)\le0$ for all $\delta\in(0,1]$, if and only if
\begin{equation}
\alpha\le1-\frac cb .
\label{eq:S-errorfree}
\end{equation}
\end{proposition}
\begin{proof}
Without errors, $\Pp CG=1$ and $\Pp DG=0$, so $r=1$ is a resident root with
$\uR=1$ and $\Delta\Phi^*=1$, and $u_{\mathrm X}(\delta)=1-\delta$.
\emph{Necessity.} By Eq.~\eqref{eq:S-near-one}, $\Rm'(1)=1-\alpha$, so
$\Lambda'(0)=1-(b/c)(1-\alpha)$, which is positive if $\alpha>1-c/b$.
\emph{Sufficiency.} With $\beta=1$ the recovery term of Eq.~\eqref{eq:S-map}
vanishes, and for $0<u\le1$
\[
\frac{\Rm_{k_1,k_2}(u)}{\Rm_{k_1,0}(u)}=\frac{u^{k_2}}{u^{k_2}+v^{k_1}(1-u^{k_2})}\le1,
\]
where subscripts denote depths. $\Rm_{k_1,0}$ is concave
(Proposition~\ref{prop:concave}) with $\Rm_{k_1,0}(1)=1$ and
$\Rm_{k_1,0}'(1)=1-\alpha$, so it lies below its tangent at $u=1$:
$\Rm_{k_1,k_2}(u)\le\Rm_{k_1,0}(u)\le1-(1-\alpha)(1-u)$. Hence
$\Lambda(\delta)=\delta-(b/c)[1-\Rm_{k_1,k_2}(1-\delta)]\le\delta[1-(b/c)(1-\alpha)]\le0$
when $\alpha\le1-c/b$.
\end{proof}
At $b/c=5$ the threshold is $\alpha=0.8$, the readout used for illustration in
Fig.~\FigOpportunism{} and Extended Data Fig.~\EDSecurity. With errors, the
resident is no longer exactly all-good, $\Delta\Phi^*<1$, and the boundary moves
(for example to $\alpha_{\mathrm U}=0.8036$ under SJ $(4,3)$). The result is
conditional on the all-good resident being selected. Stability of $r=1$ itself
requires $F'(1)=(1-\alpha)(2-p)-1<0$, that is $\alpha>0$ for SC and SS and
$\alpha>1/2$ for SH and SJ.

\subsection{Secure readout intervals}
Fix the norm, depths, $\beta$, errors and $b/c$. A readout $\alpha$ is
\emph{secure} if the high-reputation resident exists as a strictly stable root
($F=0$, $F'<0$) and satisfies $(b/c)H_{\min}\ge1$. The secure set can be empty,
can be the whole of $[0,1]$, or can consist of intervals; its boundaries are of
three kinds.
\begin{enumerate}
\item \emph{Fold (lower boundary $\alpha_{\mathrm L}$).} The high-reputation branch
ends where $F(r,\alpha)=0$ and $\partial_rF(r,\alpha)=0$, with
$F(r,\alpha)=\Rm[u_{\mathrm R}(r);\alpha,\beta,k_1,k_2]-r$, equivalently
$\Rm'[u_{\mathrm R}(r)]\,u_{\mathrm R}'(r)=1$. The fold is marginally stable and is
excluded. Because these equations contain no payoff parameter, $\alpha_{\mathrm L}$
does not depend on $b/c$. By Eq.~\eqref{eq:S-branch-sensitivity}, the resident
reputation increases with $\alpha$ along the stable branch, and the branch is
lost when $\alpha$ is lowered through $\alpha_{\mathrm L}$; in all cases computed,
the resident then falls to a low-reputation root that opportunists invade.
\item \emph{Invasion (upper boundary $\alpha_{\mathrm U}$).} $(b/c)H_{\min}=1$ on the
high branch. In every case we computed with an interior upper boundary, the
minimum was attained at $\delta\to0$, so that
\begin{equation}
F(r,\alpha_{\mathrm U})=0,\qquad
\frac bc\,\Delta\Phi(r)\,\Rm'\bigl[u_{\mathrm R}(r);\alpha_{\mathrm U}\bigr]=1 .
\label{eq:S-alphaU}
\end{equation}
For $k_2=0$ this is exact by Corollary~\ref{cor:curvature}(a). The boundary is
included (neutral mutants only).
\item \emph{Domain.} $\alpha=0$ or $\alpha=1$; included.
\end{enumerate}
Because the resident does not depend on $b/c$ and the condition
$(b/c)H_{\min}\ge1$ is monotone in $b/c$, the secure set grows with $b/c$: upper
boundaries rise with $b/c$, whereas fold boundaries stay fixed (Extended Data
Fig.~\EDSecurity b--e).

\begin{proposition}[Where the secure interval closes]\label{prop:closure}
Suppose that near a fold $(r_{\mathrm L},\alpha_{\mathrm L})$ of the high branch the
upper boundary is determined by Eq.~\eqref{eq:S-alphaU}. Then an interval of
secure readouts adjoins the fold if $b/c>u_{\mathrm R}'(r_{\mathrm L})/\Delta\Phi(r_{\mathrm L})$
and not if $b/c<u_{\mathrm R}'(r_{\mathrm L})/\Delta\Phi(r_{\mathrm L})$. Under Shunning
this critical ratio is exactly 2 for every error rate, depth and readout.
\end{proposition}
\begin{proof}
At the fold, $\Rm'u_{\mathrm R}'=1$, so the local security condition
$(b/c)\Delta\Phi\Rm'\ge1$ reads $b/c\ge u_{\mathrm R}'/\Delta\Phi$. Under SH,
$u_{\mathrm R}(r)=e_2+\kappa r^2$ and $\Delta\Phi(r)=\kappa r$
(Supplementary Table~\ref{tab:resident-functions}), so
$u_{\mathrm R}'/\Delta\Phi=2\kappa r/(\kappa r)=2$.
\end{proof}
For SJ, $u_{\mathrm R}'(r)/\Delta\Phi(r)=2-e_1(1-2e_2)/\Delta\Phi(r)$, which is
1.9710 at the fold of the selected $(4,3)$ architecture; for SC,
$u_{\mathrm R}'/\Delta\Phi=1$, so a fold never closes the interval for $b/c>1$.
These results explain why the secure intervals of SH close at $b/c=2$ in
Extended Data Fig.~\EDSecurity c, whereas that of the selected SJ architecture
is still open, though narrow, at $b/c=2$.

\subsection{Worked example: Stern Judging with depths (4,3)}
At $e_1=e_2=0.02$: $\Pp CG=0.9608$, $\Pp DG=0.02$, $\Pp CB=0.0392$ and
$\Pp DB=0.98$, so $\Delta\Phi(r)=0.9408(2r-1)$ and
$u_{\mathrm R}(r)=0.98-1.9008\,r+1.8816\,r^2$. With $\beta=1$,
\[
\Rm(u;\alpha)=\frac{u^4\sum_{m=0}^{4}[\alpha(1-u)]^m}{(1-u)^4+u^3-(1-u)^4u^3}.
\]
\begin{itemize}
\item $\alpha=1$: $\rR=0.9999996$; ALLD is excluded ($\Lambda(1)=-4.000$), but
$\Lambda(\delta)>0$ for $0<\delta<\delta_{\mathrm c}=0.365299$.
\item $\alpha=0.8$: $\rR=0.986509$ and $H_{\min}=H(0)=0.20366>c/b$; every O-DISC
mutant is excluded, with $(b/c)_{\mathrm{crit}}=1/H_{\min}=4.910$.
\item Fold: $\alpha_{\mathrm L}=0.5669858$ at $r=0.851537$. Upper boundary:
$\alpha_{\mathrm U}=0.8036040$ at $r=0.986906$, where $F'(r)<0$. The secure
interval is $0.5670<\alpha\le0.8036$.
\end{itemize}
The trade-off at $\delta=0.1$ illustrates the mechanism:

\begin{center}\small
\begin{tabular}{@{}lccccc@{}}
\toprule
Readout & $\rR$ & $r^*_{\mathrm X}(0.1)$ & Saved cost $\delta\rR$ & Lost benefit $(b/c)(\rR-r^*_{\mathrm X})$ & $\Lambda(0.1)$\\
\midrule
$\alpha=1$ & 0.9999996 & 0.999789 & 0.100000 & 0.001054 & $+0.098946$\\
$\alpha=0.8$ & 0.986509 & 0.964412 & 0.098651 & 0.110484 & $-0.011833$\\
\bottomrule
\end{tabular}
\end{center}

Supplementary Table~\ref{tab:readout} lists the secure intervals of all
architectures at $b/c=5$.

\begin{table}[htbp]
\centering\small
\caption{\textbf{Secure readout intervals at $b/c=5$} ($\beta=1$, $e_1=e_2=0.02$). Round brackets: excluded fold endpoints; square brackets: included domain or neutral-invasion endpoints. Last two columns: resident reputation and the critical benefit-to-cost ratio $1/H_{\min}$ at $\alpha=0.8$ for the selected architectures. One-step SH and SJ have lost the high-reputation resident at $\alpha=0.8$.}
\label{tab:readout}
\begin{tabular}{@{}lcccc@{}}
\toprule
Norm & One-step $(1,1)$ & Selected multi-step & $\rR(0.8)$ & $(b/c)_{\mathrm{crit}}(0.8)$\\
\midrule
SC & $(0.431929,\,0.966915]$ & $(2,0)$: $[0,\,0.807064]$ & 0.989782 & 4.827\\
SH & $(0.864863,\,0.989948]$ & $(3,0)$: $(0.591161,\,0.806464]$ & 0.986341 & 4.835\\
SJ & $(0.852446,\,0.988514]$ & $(4,3)$: $(0.566986,\,0.803604]$ & 0.986509 & 4.910\\
SS & $[0,\,0.965245]$ & $(1,2)$: $[0,\,1]$ & 0.977886 & 1.955\\
\bottomrule
\end{tabular}
\end{table}

% =====================================================================
\section{Alternative reputation filters}
\label{sec:filters}

Fig.~\FigOpportunism a compares the inertial record with two mechanisms that
also buffer isolated negative evidence. In both cases only the conversion of
assessments into a binary image changes: the norm, errors, the
private-assessment closure, the resident selection (highest root reached from
$r=1$) and Eqs.~\eqref{eq:S-odisc-u}--\eqref{eq:S-Lambda} with $\Rm$ replaced by the
alternative map are kept. The original studies use different strategy sets,
observation protocols and evolutionary dynamics, so these are controlled
adaptations, not replications; their parameters were neither optimized nor
matched in memory to the inertial record.

\subsection{Aggregation}
Following ref.~\cite{michelmata2024}, an observer samples $M$ assessments of the
donor and assigns a good image if at least $Q$ of them are positive. If each
sampled assessment is independently positive with probability $u$, the number
of positive ones is binomial, and
\begin{equation}
\Rm_{\mathrm{agg}}(u;M,Q)=\sum_{j=Q}^{M}\binom Mj u^j(1-u)^{M-j},
\qquad
\Rm_{\mathrm{agg}}'(u)=M\binom{M-1}{Q-1}u^{Q-1}(1-u)^{M-Q}.
\label{eq:S-agg}
\end{equation}
We use $(M,Q)=(2,1)$, the rule that forgives one of two negative assessments:
$\Rm_{\mathrm{agg}}(u)=1-(1-u)^2=2u-u^2$. Then
$\rR-r^*_{\mathrm X}=2(1-\uR)x+x^2$ with $x=\delta\rR\Delta\Phi^*$, and
\begin{equation}
\Lambda_{\mathrm{agg}}(\delta)=\rR\,\delta\Bigl[1-2\frac bc(1-\uR)\Delta\Phi^*
-\frac bc\,\rR\,(\Delta\Phi^*)^2\,\delta\Bigr].
\label{eq:S-agg-Lambda}
\end{equation}
For SJ at the reference parameters, $\rR=0.998188$, $\uR=0.957431$ and
$\Delta\Phi^*=0.937390$, giving
$\Lambda_{\mathrm{agg}}=0.599876\,\delta-4.377595\,\delta^2$ and invasion for
$0<\delta<0.137033$. In the error-free all-good limit,
$\Lambda_{\mathrm{agg}}=\delta[1-(b/c)\delta]$ and the interval ends at $\delta=c/b$.

\subsection{Quantitative score}
Following ref.~\cite{schmid2023}, an observer keeps an integer score
$z\in\{-z_{\max},\dots,z_{\max}\}$ for each target, increases it by one after a
positive and decreases it by one after a negative assessment (attempts to leave
the range leave the score unchanged), and reads the target as good if
$z\ge z_0$. For i.i.d.\ input $u\in(0,1)$, the score is a reflecting birth--death
chain on $j=z+z_{\max}\in\{0,\dots,2z_{\max}\}$. Detailed balance,
$\pi_j\,u=\pi_{j+1}(1-u)$, gives $\pi_j\propto t^j$ with $t=u/(1-u)$, and
\begin{equation}
\Rm_{\mathrm{score}}(u;z_{\max},z_0)=
\frac{\sum_{j=z_{\max}+z_0}^{2z_{\max}}t^j}{\sum_{j=0}^{2z_{\max}}t^j},
\qquad \Rm_{\mathrm{score}}(0)=0,\ \Rm_{\mathrm{score}}(1)=1 .
\label{eq:S-score}
\end{equation}
At $u=1/2$ the score is uniform and $\Rm_{\mathrm{score}}=(z_{\max}-z_0+1)/(2z_{\max}+1)$;
for $z_{\max}=1$, $z_0=0$, $\Rm_{\mathrm{score}}=u/(1-u+u^2)$. We use $z_{\max}=5$,
$z_0=0$. Under SJ at the reference parameters the selected resident has
$1-\rR\approx4.6\times10^{-9}$, ALLD is excluded, and O-DISC invades for
$0<\delta<0.383206$. The source model starts scores at zero; selecting the root
reached from $r=1$ is a convention of this adaptation, and SC and SH also have
lower roots for $z_{\max}=5$.

Supplementary Table~\ref{tab:fig5} summarizes the five curves of
Fig.~\FigOpportunism a.

\begin{table}[htbp]
\centering\small
\caption{\textbf{Opportunistic invasion under SJ} ($b/c=5$, $e_1=e_2=0.02$, $\beta=1$). $\delta_{\mathrm c}$: upper end of the interval with $\Lambda>0$ (0 if none).}
\label{tab:fig5}
\begin{tabular}{@{}lccc@{}}
\toprule
Mechanism & $\rR$ & $\Lambda(1)$ & $\delta_{\mathrm c}$\\
\midrule
One-step inertia $(1,1)$, $\alpha=1$ & 0.995307 & $-3.975$ & 0.008397\\
Selected $(4,3)$, $\alpha=1$ & 0.9999996 & $-4.000$ & 0.365299\\
Selected $(4,3)$, $\alpha=0.8$ & 0.986509 & $-3.946$ & 0\\
Aggregation, $(M,Q)=(2,1)$ & 0.998188 & $-3.778$ & 0.137033\\
Quantitative score, $z_{\max}=5$, $z_0=0$ & 0.99999999 & $-4.000$ & 0.383206\\
\bottomrule
\end{tabular}
\end{table}

% =====================================================================
\section{Finite-population reputation simulations}
\label{sec:simulations}

\subsection{Purpose and scope}
The simulations test how reputations form when the mean-field closure is not
imposed. Strategy counts are fixed throughout every run: the simulations
concern reputation dynamics only and contain no strategy replacement, mutation
or imitation. Strategy evolution is analysed with the replicator dynamics of
Supplementary Note~\ref{sec:strategy}. The simulations retain individual
latent records, sampling fluctuations, differences among targets of the same
strategy and correlations between the records of different observers.

\subsection{Event algorithm}
The population has $N$ individuals with fixed counts $n_{\ALLC}$, $n_{\ALLD}$ and
$n_{\DISC}$. Each DISC individual $i$ stores a latent record $L_{ij}$ for every
individual $j$, including itself; all records start in $\G$. Records of ALLC and
ALLD individuals are not stored: these strategies never consult reputations, and
observers are drawn from the DISC individuals, so omitting them changes no
event. Each event consists of the following steps.
\begin{enumerate}
\item Draw a donor uniformly from the $N$ individuals and a recipient uniformly
from the other $N-1$. Draw an observer uniformly from the $n_{\DISC}$ DISC
individuals, independently of donor and recipient; the observer may coincide
with either.
\item Determine the intended action: $\C$ for ALLC, $\D$ for ALLD; a DISC donor
samples an image from its record of the recipient,
$\mathrm{Bernoulli}[\sigma(L_{\text{donor},\text{recipient}})]$, and intends $\C$
if the image is $\G$.
\item If the intended action is $\C$, replace it by $\D$ with probability $e_1$.
\item The observer samples its own image of the recipient,
$\mathrm{Bernoulli}[\sigma(L_{\text{observer},\text{recipient}})]$, independently of
step 2, and draws a prescribed assessment that is positive with probability
$\nu_{AY}$ for the realized action $A$ and its image $Y$. The assessment is
reversed with probability $e_2$.
\item Apply $T^{+}$ or $T^{-}$ (Supplementary Table~\ref{tab:transitions}) to
$L_{\text{observer},\text{donor}}$. No other record changes.
\end{enumerate}
One sweep consists of $n_{\DISC}N$ events, the number of stored records, so each
record is updated once per sweep on average. When the DISC donor is itself the
observer, its two consultations in steps 2 and 4 are independent draws from the
same record, and the update occurs after both. The instantaneous panels of
Fig.~\FigRecognition a--d use the one-step state space with $\alpha=\beta=0$,
which generates the same image process as $(k_1,k_2)=(0,0)$
(Lemma~\ref{lem:instantaneous}).

\subsection{Estimators}
At a sampled time, let $s_{ij}=\sigma(L_{ij})$, the probability that observer $i$
currently reads target $j$ as good. Using these probabilities, rather than
sampled images, removes readout noise from the estimates. For target $j$, the
mean goodness and the probability that two distinct DISC observers, consulting
independently, read the target differently are
\begin{equation}
g_j=\frac{1}{n_{\DISC}}\sum_is_{ij},\qquad
d_j=\binom{n_{\DISC}}{2}^{-1}\sum_{i<i'}\bigl[s_{ij}(1-s_{i'j})+(1-s_{ij})s_{i'j}\bigr]
=\frac{(n_{\DISC}-1)\mu_{1,j}-\mu_{1,j}^2+\mu_{2,j}}{\binom{n_{\DISC}}{2}},
\label{eq:S-estimators}
\end{equation}
where $\mu_{1,j}=\sum_is_{ij}$ and $\mu_{2,j}=\sum_is_{ij}^2$, and $i,i'$ range over
the DISC observers (including $j$ itself if $j$ uses DISC). The second form
follows from $\sum_{i<i'}(s_i+s_{i'}-2s_is_{i'})=(n_{\DISC}-1)\mu_1-(\mu_1^2-\mu_2)$.
For deterministic readouts with a fraction $x_j$ of observers reading $j$ as
good, $d_j=2n_{\DISC}x_j(1-x_j)/(n_{\DISC}-1)$, the finite-observer version of
$2x_j(1-x_j)$. The strategy-level quantities $g_I$ and $d_I$ average $g_j$ and
$d_j$ over the $n_I$ targets of strategy $I$ and over sampling times, and
\begin{equation}
\Delta g=\frac{g_{\DISC}+g_{\ALLC}}{2}-g_{\ALLD},\qquad
\bar d=\sum_I\frac{n_I}{N}d_I .
\label{eq:S-recognition}
\end{equation}
$\Delta g$ measures how differently cooperative strategies and unconditional
defectors are treated; it is not a classification accuracy. Low $\bar d$ alone is
not evidence of good reputation formation, because observers can agree that
everyone is bad (one-step SH in Fig.~\FigRecognition b). The two statistics are
therefore interpreted jointly. $g_I$ is the finite-population counterpart of
$r_I$, but it is estimated without the closure.

\subsection{Runs, replication and uncertainty}
\begin{itemize}
\item \emph{Fig.~\FigRecognition a--d (image matrices).} $N=90$ (30 per strategy);
1,000 sweeps; seed 2027 for each panel; the matrix shows the final readout
probabilities $s_{ij}$ of the DISC observers (rows) for all targets (columns,
grouped by strategy), which are 0 or 1 at the displayed readouts. These are
illustrative single realizations.
\item \emph{Fig.~\FigRecognition e--h (depth maps).} $N=150$ (50 per strategy);
$k_1,k_2\in\{0,\dots,6\}$, $\alpha=\beta=1$; one run per condition (base seed 2027);
1,000 sweeps, sampled every 5 sweeps from sweep 800 to 1,000 inclusive (41
samples). Colour shows $\Delta g$ and circle area $\bar d$.
\item \emph{Extended Data Fig.~\EDRecognition.} The same 196 conditions with 30
independent runs each (5,880 runs). Condition seeds are generated by NumPy
\texttt{SeedSequence} from the base seed 20260918, the norm index (SC, SH, SJ,
SS), $k_1$, $k_2$ and the run index. Point estimates are means of the 30 run
means; pointwise 95\% percentile intervals use 2,000 bootstrap resamples of run
means (seed 2026091801). Samples within a run are never treated as independent.
\item \emph{Drift criterion.} For each run, the difference between the means of
the second and first halves of its sampling window is computed for $\Delta g$
and $\bar d$, and these differences are bootstrapped across runs. A condition
passes if both 95\% intervals lie inside $(-0.01,0.01)$. Conditions that failed
were rerun for 3,000, 10,000 and then 50,000 sweeps, with the first 80\% of each
run discarded and the five-sweep interval retained, until they passed. Only SJ
$(4,5)$ (3,000 sweeps) and SJ $(3,5)$ (50,000 sweeps) required extension, and
all 196 conditions passed. This is a practical check of the sampling window,
not a proof of mixing.
\end{itemize}

\subsection{Selection of representative architectures}
Among conditions that pass the drift criterion and satisfy $\Delta g\ge0.8$ and
$\bar d\le0.1$, we choose the smallest total depth $K$, then the larger
$\Delta g$, the smaller $\bar d$, the smaller $k_1$ and the smaller $k_2$. This
selects SC $(2,0)$, SH $(3,0)$, SJ $(4,3)$ and SS $(1,2)$
(Supplementary Table~\ref{tab:selected}). The cutoffs are operational criteria,
not theoretical thresholds, and selection uses no evolutionary or invasion
outcome. The stars in Fig.~\FigRecognition e--h apply the same cutoffs and
tie-breaking rule to the single-run maps (without the drift criterion) and
select the same four pairs. Extended Data Fig.~\EDRecognition e--h repeats the
selection for minimum separations $0.70,0.75,\dots,0.90$ and maximum
disagreements $0.050,0.075,\dots,0.150$.

\begin{table}[htbp]
\centering\small
\caption{\textbf{Selected architectures} (30 runs, $N=150$; brackets: pointwise 95\% bootstrap intervals).}
\label{tab:selected}
\begin{tabular}{@{}llcc@{}}
\toprule
Norm & $(k_1,k_2)$ & $\Delta g$ & $\bar d$\\
\midrule
SC & $(2,0)$ & 0.9195 [0.9192, 0.9198] & 0.0644 [0.0641, 0.0646]\\
SH & $(3,0)$ & 0.9059 [0.9056, 0.9062] & 0.0694 [0.0691, 0.0696]\\
SJ & $(4,3)$ & 0.8806 [0.8801, 0.8812] & 0.0812 [0.0809, 0.0815]\\
SS & $(1,2)$ & 0.9174 [0.9171, 0.9177] & 0.0562 [0.0560, 0.0564]\\
\bottomrule
\end{tabular}
\end{table}

\subsection{Relation to the mean-field theory}
The simulations and the theory share the event rules, records and readout,
but differ in the closure, in population size and in that strategies do not
evolve. The simulations are used to identify architectures that produce
discriminating and consistent private reputations; the strategy dynamics
then use the closed theory with these architectures. No quantitative agreement
between simulated $g_I$ and mean-field $r_I$ is claimed or used.

% =====================================================================
\section{Numerical methods}
\label{sec:numerics}

\subsection{Reputation map and roots}
$\Rm$ is evaluated from Eq.~\eqref{eq:S-map}, with arguments clipped to $[0,1]$
and the special case $\alpha=\beta=0$ evaluated as $\Rm(u)=u$; $\Rm'$ is evaluated
from Eq.~\eqref{eq:S-derivative} or by central differences. The closed form was
checked against the stationary vector of $T(u)$ computed by linear algebra for
random depths, readouts and inputs. Homogeneous-DISC roots are found by
evaluating $F$ on a uniform grid (2,001 points for the drift curves and the
depth planes, 401 points for the $161\times161$ readout plane, 4,001 points in
Fig.~\FigOpportunism), refining every sign change by bracketed root finding and
classifying roots as stable or unstable when $|F'|>10^{-7}$ and as near-fold
otherwise. In Extended Data Fig.~\EDSecurity, candidate roots are the real roots
of the polynomial numerator of $F$ (from a Chebyshev fit) and are refined by
Newton's method.

\subsection{Reputations at mixed compositions}
Eq.~\eqref{eq:S-relaxation} is integrated from $r_I=1$ by fourth-order
Runge--Kutta steps, with every stage clipped to $[0,1]$. The strategy engine
uses step 0.25 until $\max_I|\Rm(u_I)-r_I|\le10^{-11}$ (fallback step 0.1). The
simplex fields of Fig.~\FigStrategy{} use steps 0.1 (up to $2\times10^4$ steps)
and then 0.2 (up to $8\times10^4$ and a further $2\times10^5$ steps if needed),
with residual $\le10^{-12}$.

\subsection{Residents for local invasion}
The scalar flow $\mathrm dr/\mathrm d\tau=F(r)$ is integrated from $r=1$ by RK4
with step 0.18, to residual $10^{-11}$ (Extended Data Figs.~\EDRatio,
\EDDepth{} and \EDReadout i--l) or $10^{-9}$ (Extended Data Fig.~\EDErrors), and
the root is then refined by Newton steps of at most $10^{-6}$ to residual
$\le10^{-13}$. The step bound keeps the branch reached by relaxation. This
refinement matters where the resident reached from all-good is the all-bad
state ($e_2=0$ under SC and SH), at which ALLD is exactly neutral. Local outcomes
are classified from $\lambda_{\ALLD}$ and $\lambda_{\ALLC}$ (with $c=1$): ALLD is
not excluded if $\lambda_{\ALLD}\ge-10^{-10}$ (this class includes neutral
ALLD); otherwise ALLC invades if $\lambda_{\ALLC}>10^{-10}$; otherwise the
residents resist both. Dots in Extended Data Fig.~\EDDepth{} mark the last class,
and solid segments in Extended Data Fig.~\EDRatio{} mark
$R_{\mathrm L}<b/c<R_{\mathrm U}$.

\subsection{Strategy trajectories}
The engine tabulates the all-good-selected reputations exactly at the 7,381
nodes of a triangular lattice with 120 subdivisions, node $(i,j)$ being
$(f_{\ALLC},f_{\ALLD})=(i,j)/120$. A cell is flagged if any reputation component
varies by more than 0.02 among its three nodes. In unflagged cells reputations
are interpolated barycentrically; in flagged cells, which include every cell
crossed by a fold or a steep transition, they are computed exactly by relaxation
at the requested composition. Payoffs and the replicator field are computed from
these reputations. Earlier pipelines interpolated the replicator field itself
on a coarse mesh across folds, which biased sampled cooperation; this was
corrected. Trajectories are integrated by the explicit midpoint rule with
$\Delta t=0.2$, renormalizing $\boldsymbol f$ to the simplex after each stage. After
every block of 500 steps (100 time units), $\rho_{\mathrm{int}}$ is compared with
its value after the previous block; integration stops after three consecutive
changes of at most $10^{-7}$, or after 400 blocks ($4\times10^4$ time units).

\emph{Validation.} (i) Refined settings (160 subdivisions, flag threshold 0.01,
$\Delta t=0.1$, 1,000-step blocks, tolerance $10^{-8}$) changed $\bar\rho$ by at most
$1.3\times10^{-5}$ and no outcome class. (ii) An independent implementation that
recomputes reputations exactly at every step, without a lattice, reproduced the
outcome of every tested trajectory. (iii) The analytic SC no-inertia case is
reproduced. (iv) The time cap was reached only on the neutral segment of SC
without inertia (Proposition~\ref{prop:sc-line}) and by a few trajectories at
$b/c$ close to $R_{\mathrm L}$ or $R_{\mathrm U}$, where an invasion payoff is close to
zero; the outcome-gap check (Supplementary Note~\ref{sec:strategy}) holds for
every condition.

\subsection{Initial compositions}
\emph{Shifted Halton} (Fig.~\FigDepth; Extended Data Figs.~\EDResident e,
\EDDepth, \EDReadout e--h): points $i=1,\dots,512$ of the Halton sequence in bases
2 and 3, shifted modulo 1 by a uniform random vector drawn with seed 2026.
\emph{Scrambled Sobol} (Extended Data Fig.~\EDRatio): eight independently
scrambled two-dimensional Sobol sequences of 512 points (SciPy
\texttt{qmc.Sobol}~\cite{virtanen2020}, seeds 2026--2033). A point $(x,y)$ of the
unit square with $x+y>1$ is reflected to $(1-x,1-y)$, which maps the square
two-to-one onto the triangle with constant density, and
$(f_{\ALLC},f_{\ALLD},f_{\DISC})=(x,y,1-x-y)$. The same compositions are used for
all conditions within a figure.

\subsection{Opportunistic invasion}
Fig.~\FigOpportunism{} evaluates $\Lambda(\delta)$ at 5,015 distinct values of
$\delta$: 4,001 uniform points on $[0,1]$, 1,001 on $[0,0.025]$ and 100
logarithmically spaced points on $[10^{-9},10^{-3}]$. Nonzero sign changes are
refined by Brent's method (tolerance $2\times10^{-13}$). The readout scan uses
1,001 uniform values of $\alpha$ plus points next to the two boundaries (1,006
values). $\alpha_{\mathrm L}$ is obtained by solving $F=\partial_rF=0$ and
$\alpha_{\mathrm U}$ by Brent's method on Eq.~\eqref{eq:S-alphaU}; the full
$\delta$-grid confirmed that no finite $\delta$ invades inside the interval.
Residents of the alternative filters are obtained from $r=1$ by
$r\leftarrow0.8r+0.2\,\Rm_{\mathrm X}[u_{\mathrm R}(r)]$ until successive values differ
by less than $2\times10^{-14}$.

Extended Data Fig.~\EDSecurity{} uses at least 601 values of $\alpha$, refined
near folds, 161 values of $b/c\in[2,8]$ and $\beta=1$. Because $\Rm$ is rational
in $u$ and $u_{\mathrm X}$ is linear in $\delta$, $H(\delta)$ is a rational function
of $\delta$. Its minimum is found by comparing $H(0)$, $H(1)$ and $H$ at every
real root in $(10^{-8},1)$ of the numerator of $\mathrm dH/\mathrm d\delta$
(roots accepted if their imaginary part is below $10^{-7}$). A readout is
secure if $1-(b/c)H_{\min}\le10^{-10}$. This is a finite-precision search, not a
proof that no extremely narrow interior extremum exists; in all reported
cases the minimum was at $\delta\to0$.

\subsection{Software}
Analyses use Python 3.11 or later with NumPy, SciPy, pandas and Matplotlib, a
C99 implementation of the strategy engine, and Numba for the finite-population
simulations and simplex fields. The verification was run with Python 3.13.17,
NumPy 2.5.3, pandas 3.0.5, SciPy 1.18.1, Matplotlib 3.11.2 and GCC 13.3. Every
notebook loads the supplied data by default and regenerates it when its
recompute option is set.

% =====================================================================
\section{Figure-by-figure settings}
\label{sec:figures}

Supplementary Table~\ref{tab:figures} lists, for every main-text and Extended
Data figure, the model settings, the quantity shown, the numerical procedure and
the code. Unless stated, $b/c=5$ and $e_1=e_2=0.02$. Paths refer to the verified
computational package (\texttt{verified/}); working numbers of Extended Data
figures in that package differ from the manuscript numbering and are given in
parentheses.

{\small
\begin{longtable}{@{}>{\raggedright\arraybackslash}p{0.12\linewidth}>{\raggedright\arraybackslash}p{0.24\linewidth}>{\raggedright\arraybackslash}p{0.32\linewidth}>{\raggedright\arraybackslash}p{0.24\linewidth}@{}}
\caption{\textbf{Settings and procedures for each figure.}}\label{tab:figures}\\
\toprule
Figure & Model and quantity & Procedure and sample & Code\\
\midrule
\endfirsthead
\toprule
Figure & Model and quantity & Procedure and sample & Code\\
\midrule
\endhead
\bottomrule
\endfoot
Fig.~\FigModel & Schematic of instantaneous and inertial updating and readout. & Conceptual; no data. & --\\
Fig.~\FigRecognition a--d & Private images $s_{ij}$ of DISC observers; no inertia (one-step space, $\alpha=\beta=0$), full one-step and selected multi-step inertia. & Finite simulation, $N=90$ (30 per strategy), 1,000 sweeps, seed 2027; final state of one run. & \texttt{figure2\_ed1/ Figure2.ipynb}\\
Fig.~\FigRecognition e--h & $\Delta g$ (colour) and $\bar d$ (circle area) for $0\le k_1,k_2\le6$, $\alpha=\beta=1$. & $N=150$, one run per condition (base seed 2027), samples every 5 sweeps from 800 to 1,000. Stars: selection rule without drift criterion. & \texttt{figure2\_ed1/ Figure2.ipynb}\\
Fig.~\FigStrategy & Replicator dynamics; none, one-step, selected multi-step. & Field on $84\times74$ grid clipped to the simplex; reputations relaxed from all-good at every point (residual $\le10^{-12}$); vertex, edge (501-point scans) and interior equilibria; tangent stability for edge markers; strict tolerance $10^{-7}$. & \texttt{figure3\_ed2-ed5\_ed9/\allowbreak Figure3\_Simplex\_ Dynamics.ipynb}\\
Fig.~\FigDepth & $\bar\rho$ for all 180 allocations with $K\le8$, $\alpha=\beta=1$; envelope = largest $\bar\rho$ at each $K$. & Engine; 512 shifted-Halton compositions (seed 2026). & \texttt{figure4\_ed6/ Figure4\_Total\_ Depth.ipynb}; \texttt{engine/scan.py depth}\\
Fig.~\FigOpportunism a & $\Lambda(\delta)$ under SJ: one-step; $(4,3)$ with $\alpha=1$ and $0.8$; aggregation $(2,1)$; score $z_{\max}=5$, $z_0=0$; $\beta=1$. & Highest strictly stable resident (4,001-point grid, Brent); 5,015 values of $\delta$; Brent refinement of crossings; filters' residents by damped iteration. & \texttt{figure5\_ed7/ Figure5\_Opportunistic\_ Invasion.ipynb}\\
Fig.~\FigOpportunism b & $\delta_{\mathrm c}$ and $\rR$ against $\alpha$, SJ $(4,3)$, $\beta=1$. & 1,006 values of $\alpha$; $\alpha_{\mathrm L}$ from $F=\partial_rF=0$; $\alpha_{\mathrm U}$ from Eq.~\eqref{eq:S-alphaU}, checked on the full $\delta$-grid. & as above\\
ED Fig.~\EDRecognition{} (ED1) & $\Delta g$ versus $\bar d$; selection under alternative cutoffs. & $N=150$, 30 runs per condition (base seed 20260918), bootstrap (2,000 resamples, seed 2026091801), drift criterion; extended runs for SJ $(4,5)$, $(3,5)$. & \texttt{figure2\_ed1/ Extended\_Data\_ Figure1.ipynb}\\
ED Fig.~\EDResident a--d (ED2) & Drift $F(r)$ of homogeneous DISC; none, one-step, selected. & Closed-form $F$; roots from 2,001-point grids; stability from $F'$. & \texttt{figure3\_ed2-ed5\_ed9/\allowbreak Extended\_Data\_ Figure2.ipynb}\\
ED Fig.~\EDResident e (ED2) & $\bar\rho$ from all-good versus all-bad initial reputations. & Engine, quasi-static tracking protocol; 512 shifted-Halton compositions. & \texttt{engine/initial\_ reputation.py}\\
ED Fig.~\EDPlanes{} (ED9) & Highest stable homogeneous-DISC root; hatching = two stable roots. & a--d: $(1,1)$, $161\times161$ grid of $(\alpha,\beta)$, 401-point root grids; e--h: $\alpha=\beta=1$, $0\le k_1,k_2\le8$, 2,001-point root grids; full-state stability check of every depth-plane root. & \texttt{figure3\_ed2-ed5\_ed9/\allowbreak Extended\_Data\_ Figure2.ipynb}\\
ED Fig.~\EDEdge{} (ED8) & All self-consistent states on the ALLD--DISC edge; payoff advantage of DISC over ALLD on the all-good branch; folds. & Reduction to one equation in $\bar r$ (8,001-point grid, bisection) for 1,071 values of $f_{\DISC}\in[0.01,1]$. & \texttt{ed8\_reputation\_ geometry/ Extended\_Data\_ Figure8.ipynb}\\
ED Fig.~\EDRatio{} (ED4) & $\bar\rho$ against $b/c$; local protection segments; $R_{\mathrm L},R_{\mathrm U}$. & Engine; 45 ratios; eight scrambled-Sobol samples of 512 (seeds 2026--2033); mean $\pm$ s.e.\ across samples; refined residents. & \texttt{figure3\_ed2-ed5\_ed9/\allowbreak Extended\_Data\_ Figure4.ipynb}; \texttt{scan.py ratio}\\
ED Fig.~\EDErrors{} (ED5) & Local invasion classes over $(e_1,e_2)$; none, one-step, selected. & $241\times241$ grid on $[0,0.45]^2$; residents relaxed to $10^{-9}$, Newton-refined to $10^{-13}$; classification rule of Note~\ref{sec:numerics}. & \texttt{figure3\_ed2-ed5\_ed9/\allowbreak Extended\_Data\_ Figure5.ipynb}\\
ED Fig.~\EDDepth{} (ED6) & $\bar\rho$ for $0\le k_1,k_2\le7$, $\alpha=\beta=1$; dots: residents resist ALLD and ALLC. & Engine; 512 shifted-Halton compositions; dots from refined residents. & \texttt{figure4\_ed6/ Extended\_Data\_ Figure6.ipynb}\\
ED Fig.~\EDReadout a--d (ED3) & Replicator dynamics at $(1,1)$ with $(\alpha,\beta)=(1,0)$ and $(0,1)$. & As Fig.~\FigStrategy. & \texttt{Figure3\_Simplex\_ Dynamics.ipynb}\\
ED Fig.~\EDReadout e--h (ED3) & $\bar\rho$ over $41\times41$ grid of $(\alpha,\beta)$ at $(1,1)$. & Engine; 512 shifted-Halton compositions. & \texttt{scan.py readout}\\
ED Fig.~\EDReadout i--l (ED3) & Local invasion classes over $801\times801$ grid of $(\alpha,\beta)$ at $(1,1)$. & Residents relaxed to $10^{-11}$ (polynomial fallback), Newton-refined to $10^{-13}$. & \texttt{Extended\_Data\_ Figure3\_Planes.ipynb}\\
ED Fig.~\EDSecurity{} (ED7) & Security against the O-DISC family across $\alpha$ and $b/c$; colour: $\rR$. & $\ge601$ values of $\alpha$ (refined near folds), 161 values of $b/c\in[2,8]$, $\beta=1$; rational-function minimization of $H$; tolerance $10^{-10}$; fold endpoints excluded. & \texttt{figure5\_ed7/ Extended\_Data\_ Figure7.ipynb}\\
\end{longtable}
}

% =====================================================================
\section{Scope and limitations}
\label{sec:limits}

\begin{enumerate}
\item \emph{Closure.} The population-level results rest on assumptions (A1)
and (A2). The single-record map $\Rm$ is exact for i.i.d.\ inputs, and the
self-consistency equations are exact stationary equations of the closed
latent-state dynamics (Proposition~\ref{prop:equivalence}), but the closure
itself is an approximation to the finite private-record process.
\item \emph{Branch selection.} Several stationary reputation states can coexist.
All strategy-dynamics results use the branch reached from good records; the
dependence on this convention is examined in Extended Data
Fig.~\EDResident e. Stability statements refer to the auxiliary flow; they agreed
with the full latent-state linearization wherever this was checked, but the
transient dynamics of the two descriptions differ.
\item \emph{Timescales.} Strategy evolution is evaluated at stationary
reputations; simultaneous reputation and strategy change is not modelled.
\item \emph{Tested mutants.} Local resistance concerns rare ALLC, ALLD and the
i.i.d.\ O-DISC family. It does not establish resistance to other strategies,
including temporally correlated or persistent opportunism, and security against
O-DISC does not imply resistance to ALLC (selected SS architecture).
\item \emph{Sampled quantities.} $\bar\rho$ and $\omega_\C$ refer to area-uniform
initial compositions and to the stated numerical stopping rules; the
uncertainty in Extended Data Fig.~\EDRatio{} describes the sampling of initial
compositions only.
\item \emph{Selection of architectures.} The four architectures are
representative minimal-depth solutions of operational recognition criteria
within $0\le k_1,k_2\le6$, not global optima. The total-depth envelope is a
maximum over the scanned allocations, not an optimization over unrestricted
mechanisms, and inertia carries no explicit cost in the model.
\item \emph{Alternative filters.} Aggregation and quantitative scoring are
represented by single, unoptimized parameter settings and are not matched in
memory; the comparison illustrates a shared vulnerability rather than ranking
mechanisms.
\end{enumerate}

\begin{thebibliography}{9}
\bibitem{okada2018}
Okada, I., Sasaki, T. \& Nakai, Y.
A solution for private assessment in indirect reciprocity using solitary observation.
\textit{J. Theor. Biol.} \textbf{455}, 7--15 (2018).
\bibitem{radzvilavicius2019}
Radzvilavicius, A. L., Stewart, A. J. \& Plotkin, J. B.
Evolution of empathetic moral evaluation.
\textit{eLife} \textbf{8}, e44269 (2019).
\bibitem{radzvilavicius2021}
Radzvilavicius, A. L., Kessinger, T. A. \& Plotkin, J. B.
Adherence to public institutions that foster cooperation.
\textit{Nat. Commun.} \textbf{12}, 3567 (2021).
\bibitem{kessinger2023}
Kessinger, T. A., Tarnita, C. E. \& Plotkin, J. B.
Evolution of norms for judging social behavior.
\textit{Proc. Natl Acad. Sci. USA} \textbf{120}, e2219480120 (2023).
\bibitem{michelmata2024}
Michel-Mata, S., Kawakatsu, M., Sartini, J., Kessinger, T. A., Plotkin, J. B. \& Tarnita, C. E.
The evolution of private reputations in information-abundant landscapes.
\textit{Nature} \textbf{634}, 883--889 (2024).
\bibitem{kawakatsu2024}
Kawakatsu, M., Kessinger, T. A. \& Plotkin, J. B.
A mechanistic model of gossip, reputations, and cooperation.
\textit{Proc. Natl Acad. Sci. USA} \textbf{121}, e2400689121 (2024).
\bibitem{schmid2023}
Schmid, L., Ekbatani, F., Hilbe, C. \& Chatterjee, K.
Quantitative assessment can stabilize indirect reciprocity under imperfect information.
\textit{Nat. Commun.} \textbf{14}, 2086 (2023).
\bibitem{virtanen2020}
Virtanen, P. \textit{et al.}
SciPy 1.0: fundamental algorithms for scientific computing in Python.
\textit{Nat. Methods} \textbf{17}, 261--272 (2020).
\end{thebibliography}
\end{document}
